\documentclass[12pt]{book}

\usepackage[margin=1in]{geometry}
\usepackage{amsmath,amssymb,amsthm}
\usepackage{graphicx}
\usepackage{hyperref}
\usepackage{listings}
\usepackage{xcolor}
\usepackage{booktabs}
\usepackage{longtable}
\usepackage{enumitem}
\usepackage{float}
\usepackage{fancyhdr}
\usepackage{titlesec}
\usepackage{epigraph}
\usepackage{csquotes}
\setlength{\headheight}{14.5pt}

\hypersetup{
  colorlinks=true,
  linkcolor=blue!60!black,
  citecolor=green!40!black,
  urlcolor=blue!60!black
}

\lstset{
  language=C,
  basicstyle=\ttfamily\small,
  keywordstyle=\color{blue!70!black}\bfseries,
  commentstyle=\color{green!50!black}\itshape,
  stringstyle=\color{red!60!black},
  breaklines=true,
  frame=single,
  backgroundcolor=\color{gray!5},
  numbers=left,
  numberstyle=\tiny\color{gray},
  xleftmargin=2em
}

\newtheorem{definition}{Definition}[chapter]
\newtheorem{theorem}{Theorem}[chapter]
\newtheorem{proposition}{Proposition}[chapter]

\pagestyle{fancy}
\fancyhead[LE,RO]{\thepage}
\fancyhead[RE]{\leftmark}
\fancyhead[LO]{\rightmark}

\title{\Huge\textbf{The Mycelial Web} \\[1em]
\Large Trust-Weighted Decentralized Governance \\
for Post-State Civilization}
\author{Tristan Stoltz \\ Luminous Dynamics}
\date{2026}

\begin{document}

\frontmatter
\maketitle

\tableofcontents

\chapter{Preface}

The nation-state is approximately four centuries old. In that span, it has organized
humanity's largest cooperative endeavors---and its most catastrophic failures. The
Westphalian model assumes territorial sovereignty, monopoly on violence, and
hierarchical administration. These assumptions are increasingly inadequate for
problems that are global, networked, and ecological in nature.

This book presents Mycelix: a decentralized governance and services platform built
on Holochain's agent-centric architecture. With 141+ zomes across 16 cluster DNAs,
Mycelix implements a \emph{fractal civic operating system}---a CivOS---that replaces
state hierarchy with mycelial networks of mutual aid, earned trust, and commons
stewardship.

The central innovation is \emph{trust-weighted governance}: participation is not
determined by citizenship, wealth, or age, but by an 8-dimensional sovereign profile measuring
epistemic integrity, thermodynamic yield, network resilience, economic velocity, civic participation, stewardship, semantic resonance, and domain competence.
This profile gates access to five progressive tiers---Observer, Participant, Citizen,
Steward, and Guardian---each with increasing governance capabilities and vote weight.

Mycelix is implemented in approximately 785,000 lines of Rust, with SDKs in
TypeScript (226K LOC), Python, and WebAssembly. It is tested by over 8,600
automated tests across all clusters, including sweettest conductor integration tests,
property-based tests, and cross-cluster bridge verification.

This book is for system architects, governance theorists, cooperative organizers,
and anyone who believes that human coordination can evolve beyond the territorial
state.\footnote{The philosophical framework underlying these design choices, including the concept of Evolving Resonant Co-creationism and the Seven Primary Harmonies, is elaborated in Stoltz (2026), \textit{Evolving Resonant Co-creationism: A Philosophical Guide to Making It Better, Infinitely}.}

\mainmatter

%% ════════════════════════════════════════════════════════════════════════════
%% PART I: FOUNDATIONS
%% ════════════════════════════════════════════════════════════════════════════

\part{Foundations}

\chapter{Why Decentralized Governance?}

\epigraph{Neither the life of an individual nor the history of a society can be
understood without understanding both.}{C.\ Wright Mills, \textit{The Sociological
Imagination} (1959)}

\section{The Crisis of Centralized Coordination}

The 21st century confronts humanity with coordination problems that centralized
institutions were never designed to solve. Climate change requires cooperation
among 195 sovereign states whose incentives are structurally misaligned. Pandemic
response requires information sharing across jurisdictions that treat data as
sovereign territory. Supply chain resilience requires trust among parties who
have no shared legal framework.

These are not merely technical problems. They are \emph{governance} problems---failures
of the mechanisms by which humans make collective decisions and enforce collective
commitments. The nation-state, optimized for territorial control and industrial
economy, is structurally incapable of governing networked, ecological, planetary
systems.

\section{Ostrom and the Commons}

Elinor Ostrom's work on common-pool resource management demonstrated that
communities can govern shared resources without either state regulation or
market privatization. Her eight design principles for long-enduring commons
institutions provide the theoretical foundation for Mycelix:

\begin{enumerate}
  \item \textbf{Clearly defined boundaries} --- Mycelix implements this through
    identity verification and the \texttt{SovereignProfile}, which defines
    who can participate and at what level.
  \item \textbf{Proportional equivalence between benefits and costs} --- the
    trust-weighted participation system ensures that governance power scales with
    demonstrated commitment (engagement dimension).
  \item \textbf{Collective-choice arrangements} --- the Governance cluster
    implements proposals, voting, and constitutional mechanisms where all
    affected parties can participate.
  \item \textbf{Monitoring} --- Holochain's agent-centric architecture provides
    inherent auditability: every agent's source chain is a tamper-evident log.
  \item \textbf{Graduated sanctions} --- reputation decay (0.998\textsuperscript{days}),
    slashing (0.5$\times$), and blacklisting (threshold $<0.05$) implement
    progressive enforcement.
  \item \textbf{Conflict-resolution mechanisms} --- the Justice domain in the Civic
    cluster provides structured dispute resolution.
  \item \textbf{Minimal recognition of rights to organize} --- Mycelix operates
    on Holochain, which requires no permission from any authority to deploy.
  \item \textbf{Nested enterprises} --- the fractal cluster architecture
    (16 clusters, each with multiple domain zomes) implements precisely this
    principle of nested, polycentric governance.
\end{enumerate}

\section{Beyond the Nation-State}

The nation-state model assumes several properties that Mycelix explicitly rejects:

\begin{description}
  \item[Territorial sovereignty] Mycelix governance is not tied to geographic
    territory. A water stewardship council in the Commons cluster governs a
    watershed, not a political boundary.
  \item[Monopoly on legitimate violence] Mycelix replaces coercion with
    reputation consequences. The only ``force'' is reduced governance
    participation for those who lose community trust.
  \item[Hierarchical administration] Mycelix is fractal: every cluster is
    self-governing, connected to others through typed bridge dispatch rather
    than bureaucratic hierarchy.
  \item[Citizenship by birth or naturalization] Mycelix membership is earned
    through demonstrated multi-dimensional civic engagement across eight sovereign dimensions.
    The \texttt{ConsciousnessTier} system replaces birthright citizenship with
    progressive, earned governance capacity.
\end{description}

\section{The Mycelial Metaphor}

Mycelium---the vegetative body of fungi---forms vast underground networks that
connect individual organisms into cooperative webs. A single mycelial network
can span hectares, sharing nutrients between trees, signaling threats, and
decomposing dead matter into fertile soil.

Mycelix adopts this metaphor deliberately:

\begin{itemize}
  \item \textbf{No central node}: Like mycelium, Mycelix has no server, no
    headquarters, no single point of failure.
  \item \textbf{Nutrient sharing}: The Finance cluster's three currencies
    (SAP, TEND, MYCEL) flow through the network like nutrients through hyphae.
  \item \textbf{Decomposition}: Demurrage and currency composting ensure that
    hoarded value returns to the commons, just as mycelium decomposes dead
    organic matter.
  \item \textbf{Signal propagation}: Cross-cluster bridge dispatch propagates
    governance signals across domain boundaries, like chemical signals propagating
    through mycelial networks.
\end{itemize}

\section{The Great Filter of Co-Creative Wisdom}

The cosmos is silent. Despite hundreds of billions of stars in our galaxy alone,
many harboring Earth-like planets, the universe presents no evidence of other
technological civilizations. The Fermi Paradox has traditionally been explained
by two categories of ``Great Filter'': either the emergence of intelligent life
is extraordinarily improbable (the filter is behind us), or technological
civilizations inevitably destroy themselves through nuclear war, engineered
pandemics, or misaligned artificial intelligence (the filter is ahead of us).
A third possibility is more subtle and more final: civilizations fail not from
spectacular self-destruction, but from \emph{coordination failure}---the inability
to develop ethical and social technologies commensurate with their computational
and physical power. Under this model, civilizations that fail the test do not
end in galaxy-illuminating explosions. They stagnate, fragment, and devolve.
They lose the capacity to distinguish signal from noise, to trust one another,
and to coordinate collective action. To the rest of the universe, they simply
go silent. The Great Silence is not the sound of explosions but the sound of
a billion civilizations failing to mature.

Humanity currently stands at the threshold of this filter. Our technological
capacity---exponentially accelerated by artificial intelligence---vastly outpaces
our collective wisdom and ethical maturity. We possess the power to reshape our
planet and our biology, but our social operating systems remain rooted in
territorial sovereignty, zero-sum competition, and hierarchical administration.
AI, deployed to optimize these legacy systems rather than to transform them,
acts as an accelerant: it collapses shared epistemic foundations through
disinformation, erodes trust through surveillance and algorithmic bias, and
paralyzes governance through coordination failure. The result is not a single
cataclysm but a slow drowning in complexity and incoherence.

Mycelix is a concrete architectural response to the Great Filter. If the failure
mode is a collapse of trust, sense-making, and coordination, then the remedy must
be an architecture that systematically generates and scales trust through
decentralized identity, verifiable claims, and wise collective governance.
The fractal cluster design, the trust-weighted participation system, and the
Epistemic Charter described in the chapters that follow are not merely technical
conveniences---they are the social technologies that the Great Filter hypothesis
identifies as the missing prerequisite for civilizational survival. A complete
treatment of this argument appears in Stoltz (2025), \textit{The Great Filter
of Co-Creative Wisdom: A Socio-Technical Solution to the Fermi Paradox}~\cite{stoltz2025filter}.

\section{Design Philosophy: Eight Harmonies}

Mycelix is guided by the Eight Harmonies framework, which provides ethical
constraints on system design:

\begin{enumerate}
  \item \textbf{Radical Transparency}: All governance actions are auditable on
    agent source chains.
  \item \textbf{Inclusive Participation}: Trust-weighted gating ensures progressive
    access rather than binary inclusion/exclusion.
  \item \textbf{Ecological Reciprocity}: The Climate and Energy clusters embed
    environmental stewardship into economic activity.
  \item \textbf{Intergenerational Care}: Constitutional mechanisms in the Governance
    cluster can encode protections for future generations.
  \item \textbf{Distributed Power}: No single agent can accumulate unchecked
    governance authority; Guardian tier requires high scores across all eight sovereign dimensions.
  \item \textbf{Creative Expression}: The Music and Education clusters support
    human flourishing beyond mere governance.
  \item \textbf{Emergent Wisdom}: Trust-weighted participation trusts that collective
    intelligence emerges from well-structured governance access.
  \item \textbf{Sacred Stillness}: The system includes mechanisms for pause,
    reflection, and deliberation before irreversible actions.
\end{enumerate}


\chapter{Holochain Architecture}

\epigraph{Integrity doesn't require consensus.}{Arthur Brock and Eric Harris-Braun,
Holochain Whitepaper (2018)}

\section{Agent-Centric, Not Data-Centric}

Blockchain architectures are data-centric: they maintain a single, globally
consistent ledger that all participants must agree upon. This design choice
entails fundamental trade-offs---high energy consumption (proof-of-work),
limited throughput (global consensus bottleneck), and a conceptual model where
data exists independently of the agents who create it.

Holochain inverts this model. Each agent maintains their own \emph{source chain}:
a hash-linked, append-only log of their actions. Data integrity is local and
immediate---an agent signs every entry they commit. There is no global ledger
and no global consensus.

\subsection{Source Chains}

An agent's source chain is a sequence of headers, each containing:

\begin{itemize}
  \item The hash of the previous header (chain linkage)
  \item The agent's public key signature
  \item The hash of the entry data (if any)
  \item A sequence number and timestamp
\end{itemize}

This structure provides tamper evidence without requiring any other agent's
participation. If an agent modifies a historical entry, the chain of hashes
breaks, and any peer performing validation will detect the tampering.

\subsection{The DHT: Shared Awareness Without Global State}

While source chains provide local integrity, the \emph{distributed hash table}
(DHT) provides shared awareness. When an agent commits an entry, a copy is
sent to a neighborhood of peers determined by the entry's hash. These peers
\emph{validate} the entry against the DNA's validation rules before storing it.

Key properties of Holochain's DHT:

\begin{itemize}
  \item \textbf{Sharded storage}: Each agent stores only a neighborhood of the
    total data, proportional to the number of agents.
  \item \textbf{Redundancy}: Each entry is stored by multiple peers (configurable),
    providing resilience against node departure.
  \item \textbf{Eventual consistency}: The DHT converges to consistency as gossip
    propagates, but never requires synchronous global agreement.
  \item \textbf{Intrinsic validation}: Every peer independently validates entries
    against the DNA's rules. There is no ``trusted validator'' role.
\end{itemize}

\section{DNA as Application Definition}

A Holochain DNA defines the rules of a particular application space. It consists of:

\begin{description}
  \item[Integrity zomes] define entry types, link types, and validation rules.
    These are the ``constitution'' of the application: immutable, deterministic,
    and enforced by every peer.
  \item[Coordinator zomes] implement application logic---CRUD operations, queries,
    and cross-zome calls. These can be updated without changing the DNA hash.
\end{description}

This separation is crucial for Mycelix. The integrity zomes define the
\emph{rules of governance}---what constitutes a valid proposal, a valid vote,
a valid identity credential. The coordinator zomes implement the
\emph{mechanisms of governance}---how proposals are submitted, how votes are
tallied, how credentials are issued.

\section{hApp Architecture}

A Holochain application (hApp) bundles one or more DNAs into roles. Each role
is instantiated as a \emph{cell} on the conductor---the Holochain runtime.
Cells within the same conductor can communicate via \texttt{call(CallTargetCell, ...)}.

Mycelix uses a \emph{unified hApp} architecture: all 16 cluster DNAs are
bundled into a single hApp with defined roles. This enables cross-cluster
communication via \texttt{CallTargetCell::OtherRole}, avoiding the complexity
of inter-hApp signaling.

\begin{lstlisting}[caption={Unified hApp manifest (simplified)},language={}]
---
manifest_version: "1"
name: mycelix-unified
roles:
  - name: commons
    dna: mycelix-commons.dna
  - name: civic
    dna: mycelix-civic.dna
  - name: identity
    dna: mycelix-identity.dna
  - name: governance
    dna: mycelix-governance.dna
  - name: finance
    dna: mycelix-finance.dna
  # ... 11 more cluster roles
\end{lstlisting}

\section{Why Not Blockchain?}

The choice of Holochain over blockchain (Ethereum, Solana, etc.) is deliberate
and follows from Mycelix's design requirements:

\begin{longtable}{p{4cm}p{4.5cm}p{4.5cm}}
\toprule
\textbf{Property} & \textbf{Blockchain} & \textbf{Holochain} \\
\midrule
\endhead
Consensus & Global (all nodes agree) & Local (source chain + DHT validation) \\
Scaling & O(n) per transaction & O(1) per agent action \\
Data model & Global state (ledger) & Agent-centric (source chains) \\
Energy & PoW: enormous; PoS: moderate & Minimal (no mining) \\
Governance & Token-weighted (plutocratic) & Identity-weighted (meritocratic) \\
Privacy & Pseudonymous but public & Private source chains + selective sharing \\
Offline & Impossible (requires consensus) & Possible (local source chain) \\
\bottomrule
\end{longtable}

The most critical difference for governance applications is the consensus
model. Blockchain governance is inherently plutocratic: those with more tokens
have more votes. Holochain's agent-centric model allows Mycelix to implement
trust-weighted governance where influence derives from identity, reputation,
community trust, and engagement rather than wealth.


\chapter{The Fractal Cluster Design}

\epigraph{A complex system that works is invariably found to have evolved from
a simple system that worked.}{John Gall, \textit{Systemantics} (1975)}

\section{Why 16 Clusters?}

Mycelix organizes its 133+ zomes into 16 cluster DNAs. This is neither arbitrary
nor aspirational---it emerges from two constraints:

\begin{enumerate}
  \item \textbf{Holochain's cell model}: Each DNA is instantiated as a cell.
    Cross-zome calls within a cell (\texttt{CallTargetCell::Local}) are fast
    and synchronous. Cross-cell calls (\texttt{CallTargetCell::OtherRole})
    require message passing.
  \item \textbf{Domain cohesion}: Zomes that frequently interact should share
    a DNA. A housing governance vote that checks property records should not
    require a cross-cluster call.
\end{enumerate}

The 16 clusters group by domain affinity:

\begin{longtable}{llrl}
\toprule
\textbf{Cluster} & \textbf{Domain} & \textbf{Zomes} & \textbf{Status} \\
\midrule
\endhead
mycelix-commons & Property, Housing, Care, Mutual Aid, Water, Food, Transport & 39 & Built \\
mycelix-civic & Justice, Emergency, Media, Resonance Feed & 18 & Built \\
mycelix-hearth & Kinship, Gratitude, Care, Autonomy, Decisions & 12 & Built \\
mycelix-identity & DID, MFA, Trust Credentials, Recovery, Web-of-Trust & 13 & Built \\
mycelix-governance & Proposals, Voting, DKG, Councils, Constitution & 7 & Built \\
mycelix-finance & SAP/TEND/MYCEL Payments, Treasury, Staking & 8 & Built \\
mycelix-health & FHIR, CDS, Telehealth, Trials, Insurance & 15 & Built \\
mycelix-knowledge & Claims, Graph, Query, Inference, Factcheck & 8 & Built \\
mycelix-personal & Identity Vault, Health Vault, Credential Wallet & 4 & Built \\
mycelix-attribution & Dependency Registry, Usage Receipts, Reciprocity & 3 & Built \\
mycelix-mail & PQC-Encrypted Decentralized Email & 12 & Built \\
mycelix-supplychain & Provenance Tracking & 8 & Built \\
mycelix-marketplace & Arbitration & 8 & Built \\
mycelix-energy & Projects, Investments, Grid, Regenerative & 5 & Built \\
mycelix-climate & Carbon Credits, Projects & 3 & Built \\
mycelix-space & Orbital Objects, Debris Bounties, Traffic Control & 5 & Built \\
\bottomrule
\end{longtable}

Additional clusters (Music, Education, DeSci, Core) bring specialized
capabilities for cultural production, learning, and research.

\section{Anatomy of a Cluster}

Each cluster follows a consistent internal architecture:

\begin{description}
  \item[Integrity zomes] (one per domain) define entry types and validation rules.
    For example, \texttt{housing\_integrity} defines \texttt{HousingUnit},
    \texttt{Membership}, and \texttt{MaintenanceRequest} entry types.
  \item[Coordinator zomes] (one per domain) implement CRUD operations, queries,
    and domain logic. For example, \texttt{housing\_units} provides
    \texttt{create\_unit}, \texttt{list\_units}, \texttt{transfer\_unit}.
  \item[Bridge zome] (one per cluster) handles cross-cluster dispatch. It
    receives typed requests from other clusters and routes them to the
    appropriate domain zome via local calls.
  \item[Shared types crate] (\texttt{mycelix-bridge-entry-types}) defines DHT
    entry types shared across clusters.
\end{description}

\section{The Fractal Pattern}

The architecture is fractal in a precise sense: the same governance pattern
repeats at every scale.

\begin{itemize}
  \item \textbf{Zome level}: Each coordinator zome checks the caller's
    trust credential before performing governance-sensitive operations.
  \item \textbf{Cluster level}: Each cluster's bridge zome enforces cluster-wide
    policies (rate limiting, allowlists) before routing requests.
  \item \textbf{System level}: The unified hApp manifest defines which clusters
    can communicate, and the bridge-common crate provides type-safe dispatch
    across all clusters.
\end{itemize}

This fractal structure means that a water stewardship circle, a city-wide
housing cooperative, and the global Mycelix network all use the same governance
primitives: trust profiles, tiered access, typed bridge dispatch.

\section{Shared Infrastructure: bridge-common}

The \texttt{crates/mycelix-bridge-common} crate (295+ tests) provides the shared
infrastructure that all clusters depend on:

\begin{itemize}
  \item \texttt{SovereignProfile} --- the 8D sovereign profile (epistemic integrity, thermodynamic yield, network resilience, economic velocity, civic participation, stewardship, semantic resonance, domain competence)
    profile (detailed in Chapter 4)
  \item \texttt{ConsciousnessTier} --- the 5-tier enum (Observer through Guardian)
  \item \texttt{ConsciousnessCredential} --- time-limited, BLAKE3-committed credentials
  \item \texttt{BridgeDomain} --- typed domain routing (12 variants)
  \item \texttt{CommonsZome} / \texttt{CivicZome} --- typed zome routing (38 + 15 variants)
  \item Rate limiting, allowlists, and reputation primitives
\end{itemize}


%% ════════════════════════════════════════════════════════════════════════════
%% PART II: TRUST-WEIGHTED GOVERNANCE
%% ════════════════════════════════════════════════════════════════════════════

\part{Trust-Weighted Governance}

\chapter{The 8D Sovereign Profile}
\label{ch:trust-profile}

\epigraph{Who should govern? is the wrong question. The right question is:
How can we organize institutions so that bad or incompetent rulers can be
prevented from doing too much damage?}{Karl Popper, \textit{The Open Society
and Its Enemies} (1945)}

\section{Terminology and Scope}

The codebase uses \texttt{SovereignProfile} and \texttt{CivicTier} for the types described in this chapter, replacing the earlier \texttt{ConsciousnessProfile} (4D) with an eight-dimensional model. The eight dimensions are each measured directly from a primary source cluster, not from self-reports or financial proxies:

\begin{itemize}[nosep]
  \item $D_0$: \textbf{Epistemic Integrity} --- truth-telling track record (knowledge cluster)
  \item $D_1$: \textbf{Thermodynamic Yield} --- physical energy contribution (energy/grid)
  \item $D_2$: \textbf{Network Resilience} --- node uptime and bandwidth (mesh-time)
  \item $D_3$: \textbf{Economic Velocity} --- anti-hoarding TEND circulation (finance)
  \item $D_4$: \textbf{Civic Participation} --- voting, jury duty, proposals (governance)
  \item $D_5$: \textbf{Stewardship \& Care} --- verified commons labor (attribution)
  \item $D_6$: \textbf{Semantic Resonance} --- federated learning Hamming similarity (core-FL)
  \item $D_7$: \textbf{Domain Competence} --- peer-verified expertise with Ebbinghaus decay (craft)
\end{itemize}

These names reflect Mycelix's architecture as a consciousness-coupled governance platform. In the standalone Mycelix context---where no Symthaea engine is connected---these eight dimensions function as \emph{earned trust-weighted participation gates}. The mechanism is straightforward: agents accumulate scores across eight physically-grounded dimensions, and the composite determines governance capability. The \texttt{gate\_civic()} function provides backward-compatible gating: it attempts native 8D evaluation first, falling back to the legacy 4D \texttt{ConsciousnessProfile} with automatic conversion.

A reader who is skeptical of computational consciousness can safely read the sovereign profile as a ``multi-dimensional earned trust capacity profile.'' The eight dimensions, the tier system, the sigmoid vote weighting, and the credential mechanism are all defined by their engineering properties, not by any metaphysical commitment. The system's value lies in its measurable governance outcomes: Sybil resistance, progressive access, reputation accountability, and community-weighted trust. These properties hold regardless of one's position on machine consciousness.

\section{Motivation: Beyond Binary Access Control}

Traditional access control is binary: you either have permission or you do not.
Role-based access control (RBAC) adds granularity but remains fundamentally
categorical. A user is either ``admin'' or ``editor'' or ``viewer.''

This model is inadequate for governance. Governance participation should be
\emph{progressive}---increasing with demonstrated trustworthiness, commitment,
and community standing. A newcomer should be able to observe and learn before
participating. A long-standing community member should have more governance
weight than a newcomer, but not infinite weight.

Mycelix replaces binary access control with an \emph{8-dimensional sovereign
profile} that produces a continuous governance capability score. No single dimension's weight can exceed 50\% of the composite---a constitutional invariant that prevents single-axis capture.

\section{The Eight Dimensions}

Each dimension is a continuous value in $[0, 1]$, measured from a primary source cluster:

\begin{definition}[SovereignProfile]
A trust profile $\mathbf{P} = (D_0, D_1, \ldots, D_7) \in [0,1]^8$ where:
\begin{itemize}
  \item $D_0$ (Epistemic Integrity) $\in [0,1]$: Truth-telling track record from validated claims in the knowledge cluster. Saturates at 50 validated claims.
  \item $D_1$ (Thermodynamic Yield) $\in [0,1]$: Physical energy contribution from verified energy/grid records. Saturates at 30 verified records.
  \item $D_2$ (Network Resilience) $\in [0,1]$: Node uptime and bandwidth provision from mesh-time logs. Saturates at 720 time anchors.
  \item $D_3$ (Economic Velocity) $\in [0,1]$: Anti-hoarding compliance from TEND mutual credit exchanges. Saturates at 50 TEND exchanges.
  \item $D_4$ (Civic Participation) $\in [0,1]$: Voting, jury duty, and proposal activity from governance records. Saturates at 20 governance actions.
  \item $D_5$ (Stewardship \& Care) $\in [0,1]$: Verified commons labor from attribution/reciprocity pledges. Saturates at 30 pledges.
  \item $D_6$ (Semantic Resonance) $\in [0,1]$: Community value alignment via Hamming similarity between agent and consensus training hypervectors in federated learning.
  \item $D_7$ (Domain Competence) $\in [0,1]$: Peer-verified expertise from living credentials with Ebbinghaus decay. Saturates at 3 living credentials.
\end{itemize}
\end{definition}

\paragraph{Constitutional Decay.} All dimensions decay exponentially without engagement: $S_{\text{decayed}} = S_{\text{raw}} \cdot e^{-\lambda \cdot \Delta t}$, where $\lambda$ is community-configurable within constitutional bounds ($\lambda_{\min} = 0.001$/day, half-life 693 days; $\lambda_{\max} = 0.020$/day, half-life 35 days). Lambda $= 0$ is constitutionally impossible, preventing permanent oligarchy.

\paragraph{Privacy.} The sovereign profile never leaves the participant's device. Governance proofs use zero-knowledge STARK proofs ($\sim$1ms prove, $\sim$0.3ms verify, $\sim$7.5KB proof) that demonstrate threshold compliance without revealing raw dimensions.

\paragraph{Legacy Compatibility.} The earlier 4D \texttt{ConsciousnessProfile} (Identity, Reputation, Community, Engagement) remains in the codebase for backward compatibility. The \texttt{gate\_civic()} function evaluates the native 8D profile first, falling back to 4D with automatic conversion:
$\text{Identity} \to D_0 + D_2$, $\text{Reputation} \to D_3 + D_5$, $\text{Community} \to D_4 + D_6$, $\text{Engagement} \to D_1 + D_7$.

\subsection{Identity Dimension}

The identity dimension measures how strongly an agent has verified their identity.
It maps directly from the MFA (Multi-Factor Authentication) assurance level
defined in the Identity cluster:

\begin{itemize}
  \item \textbf{Anonymous} ($I = 0.0$): No identity verification. The agent has
    a Holochain public key but nothing else.
  \item \textbf{Basic} ($I = 0.25$): Single-factor authentication (e.g., email
    verification).
  \item \textbf{Verified} ($I = 0.50$): Multi-factor authentication (e.g.,
    email + TOTP).
  \item \textbf{HighlyAssured} ($I = 0.75$): Strong multi-factor (e.g.,
    hardware security key + biometric).
  \item \textbf{Critical} ($I = 1.0$): Maximum assurance, including identity
    proofing by trusted community members.
\end{itemize}

This dimension provides Sybil resistance: creating fake identities becomes
progressively more expensive as higher assurance levels require more verification.

\subsection{Reputation Dimension}

The reputation dimension aggregates an agent's behavioral history across all
Mycelix clusters. It uses exponential decay with a 30-day half-life:

\begin{equation}
  R(t) = \sum_{i} w_i \cdot r_i \cdot \left(\frac{1}{2}\right)^{(t - t_i) / 30}
\end{equation}

where $r_i$ is the reputation event value, $t_i$ is the event timestamp (in days),
$w_i$ is the source weight, and $t$ is the current time.

Key properties:
\begin{itemize}
  \item \textbf{Temporal decay}: Reputation fades with time, preventing ancient
    good behavior from conferring permanent privilege.
  \item \textbf{Multi-source}: Reputation events come from all clusters, weighted
    by domain relevance.
  \item \textbf{Slashing}: Severe misbehavior triggers a 0.5$\times$ multiplier
    on accumulated reputation.
  \item \textbf{Blacklist threshold}: Agents whose reputation falls below 0.05
    are blacklisted from governance.
  \item \textbf{Restoration}: Blacklisted agents can restore reputation through
    100 verified positive interactions.
\end{itemize}

\subsection{Community Dimension}

The community dimension measures peer trust---how much the agent's community
vouches for them. Trust attestations are weighted by the attestor's own
governance tier:

\begin{equation}
  C = \frac{\sum_{j} \tau_j \cdot w(\text{tier}_j)}{\sum_{j} w(\text{tier}_j)}
\end{equation}

where $\tau_j$ is attestor $j$'s trust rating and $w(\text{tier}_j)$ is the
weight assigned to their tier. This creates a virtuous cycle: attestations
from higher-tier community members carry more weight, incentivizing genuine
community building rather than attestation farming.

\subsection{Engagement Dimension}

The engagement dimension measures domain-specific participation. Unlike the
other three dimensions (which are global), engagement is computed locally
by each cluster's bridge zome. This means an agent can have high engagement
in the Health cluster but low engagement in the Finance cluster.

For Symthaea-connected agents, the engagement dimension maps from the unified
consciousness score ($C_{\text{unified}}$), creating a bridge between Symthaea's
consciousness research and Mycelix governance participation.

\section{Combined Score}

The combined score is a weighted average:

\begin{equation}
  S = 0.25 \cdot I + 0.25 \cdot R + 0.30 \cdot C + 0.20 \cdot E
\end{equation}

The weights reflect a deliberate design choice: community trust (30\%) is
weighted highest because Mycelix is fundamentally a system of \emph{mutual
recognition}. Identity and reputation (25\% each) provide the foundation.
Engagement (20\%) ensures active participation matters but cannot substitute
for the other dimensions.

\section{Continuous Vote Weight}

Rather than using hard tier thresholds for vote weight, Mycelix employs a
sigmoid function centered at the Citizen threshold (0.4):

\begin{equation}
  W(S) = \frac{W_{\max}}{1 + e^{-(S - 0.4) / T}}
\end{equation}

where $W_{\max} = 10{,}000$ basis points and $T = 0.05$ is the temperature
parameter. This provides:

\begin{itemize}
  \item Smooth transitions (no ``cliff'' at tier boundaries)
  \item Noise tolerance (small score fluctuations don't cause dramatic weight changes)
  \item Configurable strictness (lower $T$ = sharper transition)
\end{itemize}


\chapter{Five Governance Tiers}

\section{Tier Definitions}

\begin{definition}[ConsciousnessTier]
Five governance tiers, derived from the combined trust score $S$:
\begin{enumerate}
  \item \textbf{Observer} ($S < 0.3$): Read-only access. Can browse public data
    and observe governance proceedings.
  \item \textbf{Participant} ($S \geq 0.3$): Can submit basic proposals and
    comment on governance actions. Vote weight: 5,000 bp (50\%).
  \item \textbf{Citizen} ($S \geq 0.4$): Full voting rights on standard
    governance actions. Vote weight: 7,500 bp (75\%).
  \item \textbf{Steward} ($S \geq 0.6$): Can initiate constitutional amendments,
    manage council membership, and access cross-cluster governance. Vote weight:
    10,000 bp (100\%).
  \item \textbf{Guardian} ($S \geq 0.8$): Emergency governance powers including
    circuit-breaker activation, emergency fund release, and threat response
    coordination. Vote weight: 10,000 bp (100\%).
\end{enumerate}
\end{definition}

\section{Hysteresis: Preventing Tier Oscillation}

A naive threshold system would cause agents near tier boundaries to oscillate
rapidly as their scores fluctuate from measurement noise. Mycelix implements
hysteresis with a margin of $\pm 0.05$:

\begin{itemize}
  \item \textbf{Promotion} requires crossing $\text{threshold} + 0.05$
  \item \textbf{Demotion} requires dropping below $\text{threshold} - 0.05$
\end{itemize}

An agent at Citizen tier (threshold 0.4) must reach 0.45 to promote to Steward,
but only drops to Citizen if they fall below 0.55 ($0.6 - 0.05$). This creates
a stability band around each tier boundary.

\section{Governance Requirements}

Each governance action specifies requirements as a \texttt{GovernanceRequirement}
struct:

\begin{lstlisting}
pub struct GovernanceRequirement {
    pub min_tier: ConsciousnessTier,
    pub min_identity: Option<f64>,
    pub min_community: Option<f64>,
    pub min_reputation: Option<f64>,
    pub min_engagement: Option<f64>,
}
\end{lstlisting}

This allows fine-grained control. A constitutional amendment might require:
\begin{itemize}
  \item Minimum tier: Steward ($S \geq 0.6$)
  \item Minimum identity: 0.75 (HighlyAssured)
  \item Minimum community: 0.5 (strong peer attestation)
\end{itemize}

\section{Credentials: Time-Limited Trust}

Trust profiles are not queried in real-time (which would require cross-cluster
calls during governance actions). Instead, each agent's bridge zome issues a
\texttt{ConsciousnessCredential}---a time-limited (24-hour default) credential
containing the profile, derived tier, and a BLAKE3 trajectory commitment.

The credential is stored on the agent's source chain and presented to governance
zomes as proof of tier. This design:

\begin{itemize}
  \item Eliminates cross-cluster calls during governance actions
  \item Provides a 24-hour validity window (stale data is bounded)
  \item Includes a BLAKE3 hash of the agent's behavioral trajectory, preventing
    credential transfer between agents
\end{itemize}

\section{Bootstrap Mode}

New communities face a bootstrapping problem: no one has reputation or community
trust yet. Mycelix handles this with bootstrap mode:

\begin{itemize}
  \item Communities with fewer agents than \texttt{BOOTSTRAP\_COMMUNITY\_THRESHOLD}
    operate under relaxed gating
  \item Minimum identity of \texttt{BOOTSTRAP\_MIN\_IDENTITY} is still required
    (Sybil resistance is never relaxed)
  \item Bootstrap credentials have a separate TTL (\texttt{BOOTSTRAP\_TTL\_US})
  \item As the community grows past the threshold, standard gating activates
    automatically
\end{itemize}


\chapter{Bridge-Common: Type-Safe Cross-Cluster Dispatch}

\section{The Problem: Fragile String Matching}

In early Mycelix prototypes, cross-cluster dispatch used string matching:

\begin{lstlisting}
if query_type.contains("transfer") {
    call(housing_transfer_zome, ...)
} else if query_type.contains("membership") {
    call(housing_membership_zome, ...)
}
\end{lstlisting}

This was fragile: a query type of \texttt{"transfer\_membership"} would match
both branches. Case sensitivity caused bugs. New zomes required updating
string patterns across multiple clusters.

\section{Typed Routing}

The \texttt{bridge-common} crate replaces string matching with strongly typed
enums:

\begin{lstlisting}
pub enum BridgeDomain {
    Property, Housing, Care, Mutualaid,
    Water, Food, Transport, Support,
    Space, Justice, Emergency, Media,
}

pub enum CommonsZome {
    PropertyRegistry, PropertyTransfer,
    PropertyDisputes, PropertyCommons,
    HousingUnits, HousingMembership,
    // ... 38 total variants
}
\end{lstlisting}

Each zome variant maps to a canonical snake\_case string matching the actual
Holochain zome name. The routing function \texttt{resolve\_commons\_zome}
performs case-insensitive domain parsing via \texttt{BridgeDomain::from\_str\_loose},
which handles variations like ``mutualaid'', ``mutual\_aid'', and ``mutual-aid''.

\section{Cross-Cluster Roles}

Communication between clusters uses Holochain's \texttt{CallTargetCell::OtherRole}
mechanism. The \texttt{CrossClusterRole} enum defines the available roles:

\begin{lstlisting}
pub enum CrossClusterRole {
    Commons,
    Civic,
    Identity,
    Governance,
    Finance,
    // ... all 16 cluster roles
}
\end{lstlisting}

Each cross-cluster call specifies the target role, target zome, and function
name. The bridge-common crate provides allowlists that restrict which
source-destination pairs are permitted, preventing unauthorized cross-cluster
access.

\section{Rate Limiting and Allowlists}

Bridge dispatch includes two security mechanisms:

\begin{description}
  \item[Allowlists] specify which source clusters can call which destination
    zomes. For example, only the Governance cluster can call the Identity
    cluster's trust credential verification functions.
  \item[Rate limiting] prevents any single agent from overwhelming a cluster's
    bridge with requests. Limits are configurable per domain and per agent tier.
\end{description}

\section{The 295+ Test Suite}

The bridge-common crate is one of the most heavily tested components in Mycelix,
with 295+ tests covering:

\begin{itemize}
  \item Profile construction and validation (NaN handling, clamping)
  \item Tier derivation and hysteresis
  \item Credential issuance and expiry
  \item Domain routing (case insensitivity, unknown domains)
  \item Vote weight computation (sigmoid properties, edge cases)
  \item Reputation decay, slashing, and blacklisting
  \item Property-based tests (tier monotonicity, score bounds)
\end{itemize}


\chapter{Collective Intelligence}

\epigraph{None of us is as smart as all of us.}{Ken Blanchard}

\section{Beyond Individual Trust}

The preceding chapters describe how individual agents earn trust through the 4D
profile. But governance is a \emph{collective} phenomenon---the quality of a
community's decisions depends not only on the trustworthiness of individual
participants but on the coherence of the group as a whole. Mycelix addresses
this through three mechanisms: collective Phi, fragile consensus detection, and
the Epistemic Charter.

\section{Collective Phi: Group Consciousness Vectors}

When Symthaea is connected, each agent produces a consciousness state
characterized by six measurable dimensions: integrated information ($\Phi$),
workspace activation (GWT), attention coherence, meta-cognitive depth,
emotional valence, and social-model accuracy. These dimensions form a
6-dimensional \emph{agent consciousness vector}:

\begin{equation}
  \mathbf{c}_i = (\Phi_i,\; W_i,\; A_i,\; M_i,\; V_i,\; S_i)
\end{equation}

For any pair of agents $i, j$, the pairwise coherence is the cosine similarity
of their consciousness vectors:

\begin{equation}
  \kappa_{ij} = \frac{\mathbf{c}_i \cdot \mathbf{c}_j}{\|\mathbf{c}_i\| \, \|\mathbf{c}_j\|}
\end{equation}

The \emph{collective Phi} of a governance quorum $Q$ is defined as:

\begin{equation}
  \Phi_Q = \frac{2}{|Q|(|Q|-1)} \sum_{i < j \in Q} \kappa_{ij}
\end{equation}

This metric quantifies how aligned a deliberating group's cognitive states are.
A collective Phi above 0.7 indicates coherent deliberation; below 0.3 suggests
fragmentation that may produce low-quality decisions.

\section{Fragile Consensus Detection}

Not all consensus is genuine. A vote that passes with 51\% in favor may reflect
deep division rather than collective agreement. Mycelix detects fragile consensus
by analyzing:

\begin{description}
  \item[Vote margin] The gap between the winning and losing tallies. Margins
    below a configurable threshold trigger extended deliberation periods.
  \item[Tier distribution] A proposal that passes only because Guardian-tier
    agents outvoted many Citizen-tier agents may reflect elite capture rather
    than community consensus.
  \item[Coherence gap] When collective Phi is low but the vote margin is high,
    the consensus may be fragile---agents voted the same way for different
    reasons and may defect when circumstances change.
  \item[Temporal stability] If pre-vote sentiment polling shows rapid opinion
    changes, the consensus may not be durable.
\end{description}

When fragile consensus is detected, the Governance cluster can automatically:
\begin{enumerate}
  \item Extend the voting period
  \item Require a higher approval threshold
  \item Flag the decision for review after a cooling-off period
  \item Recommend facilitated dialogue before re-vote
\end{enumerate}

\section{Reputation System}

The reputation system provides the temporal dimension of collective intelligence.
It ensures that trust is earned, maintained, and can be lost:

\begin{description}
  \item[Exponential decay] Reputation decays at $0.998^{\text{days}}$, yielding
    a half-life of approximately 346 days. This prevents ancient contributions
    from conferring permanent influence while allowing sustained good behavior
    to accumulate meaningful trust.
  \item[Slashing] Severe violations---data falsification, Sybil attacks,
    governance manipulation---trigger a 0.5$\times$ multiplicative slash on
    accumulated reputation. The slash is immediate and recorded on the agent's
    source chain.
  \item[Blacklisting] Agents whose reputation falls below 0.05 lose all
    governance participation. This threshold is deliberately low: it takes
    sustained or severe misbehavior to reach it.
  \item[Restoration] Blacklisted agents can restore reputation through 100
    verified positive interactions across at least two clusters. This ensures
    that restoration requires genuine community re-engagement, not merely
    waiting out a timer.
\end{description}

\section{The Epistemic Charter (E-N-M Classification)}

Governance decisions vary in their epistemic character. A proposal to fund a
bridge repair is fundamentally different from a proposal to change the community's
constitutional values. The Epistemic Charter provides a three-axis classification
for governance claims:

\begin{description}
  \item[E (Empirical)] Scored E0--E4, from Subjective (E0: personal testimony)
    to Publicly Reproducible (E4: independently verified by multiple observers).
    Higher E levels require stronger evidence and more verifiers.
  \item[N (Normative)] Scored N0--N3, from Personal preference (N0) to Axiomatic
    principles (N3: foundational values that define the community). Higher N
    levels require broader consensus and longer deliberation.
  \item[M (Materiality)] Scored M0--M3, from Ephemeral (M0: temporary, easily
    reversible) to Foundational (M3: permanent, defines the community's
    structure). Higher M levels require higher-tier approval and constitutional
    protections.
\end{description}

A governance action's E-N-M classification determines:
\begin{itemize}
  \item The minimum tier required to propose and vote
  \item The evidence standard for supporting claims
  \item The approval threshold (simple majority for low-M, supermajority for high-M)
  \item The deliberation period (hours for M0, weeks for M3)
\end{itemize}

For example, a claim that ``the water pH in the eastern watershed is 6.8'' is
classified E3-N0-M1 (empirically measurable, no normative content, moderate
materiality). A proposal that ``the community shall adopt a policy of
restorative justice'' is E0-N3-M3 (not empirically testable, deeply normative,
foundational). The governance requirements differ accordingly.


%% ════════════════════════════════════════════════════════════════════════════
%% PART III: CORE CLUSTERS
%% ════════════════════════════════════════════════════════════════════════════

\part{Core Clusters}

\chapter{Identity}

\section{Overview}

The Identity cluster is the foundation of Mycelix. Every other cluster depends
on it for agent verification, trust credentials, and the identity dimension of
the trust profile. It contains 13 zomes across 11 domains plus 1 bridge
and 1 shared:

\begin{enumerate}
  \item \textbf{did\_registry} --- DID (Decentralized Identifier) management.
    Each agent registers a \texttt{did:mycelix:<pubkey>} identifier.
  \item \textbf{mfa} --- Multi-Factor Authentication with configurable assurance
    levels (Anonymous through Critical).
  \item \textbf{trust\_credential} --- Peer-to-peer trust attestations that feed
    the community dimension.
  \item \textbf{verifiable\_credential} --- W3C Verifiable Credential issuance
    and verification.
  \item \textbf{credential\_schema} --- Schema registry for credential types.
  \item \textbf{recovery} --- Social recovery mechanism (threshold of trusted
    peers can restore lost keys).
  \item \textbf{revocation} --- Credential and key revocation with DHT propagation.
  \item \textbf{education} --- Educational credential management.
  \item \textbf{name-registry} --- Human-readable name registration and resolution.
  \item \textbf{web-of-trust} --- Transitive trust computation across the network.
  \item \textbf{bridge} --- Cross-cluster identity verification dispatch.
\end{enumerate}

\section{DID Architecture}

Mycelix uses a method-specific DID scheme: \texttt{did:mycelix:<pubkey>}.
Each DID document is stored on the agent's source chain and replicated to
the DHT. The document contains:

\begin{itemize}
  \item The agent's public key (Ed25519)
  \item Optional post-quantum key (ML-KEM-768, for future-proofing)
  \item Service endpoints (which clusters the agent participates in)
  \item Authentication methods (linked MFA assurance level)
  \item Delegation chains (sub-passports for authorized devices)
\end{itemize}

\section{Multi-Factor Authentication}

The MFA zome implements five assurance levels, each requiring progressively
stronger verification:

\begin{longtable}{lll}
\toprule
\textbf{Level} & \textbf{Identity Score} & \textbf{Requirements} \\
\midrule
Anonymous & 0.00 & Holochain key only \\
Basic & 0.25 & Email or phone verification \\
Verified & 0.50 & Two independent factors \\
HighlyAssured & 0.75 & Hardware key + biometric \\
Critical & 1.00 & In-person proofing by trusted peers \\
\bottomrule
\end{longtable}

\section{Web of Trust}

The web-of-trust zome computes transitive trust across the network. If Alice
trusts Bob and Bob trusts Carol, Alice has indirect trust in Carol, discounted
by a configurable decay factor per hop. This enables:

\begin{itemize}
  \item Discovery of trustworthy agents outside one's direct community
  \item Sybil detection (clusters of mutually-attesting fake identities become
    visible as isolated subgraphs)
  \item Community dimension computation (aggregated trust from the web)
\end{itemize}

\section{Social Recovery}

The recovery zome implements $m$-of-$n$ social recovery: an agent designates
$n$ trusted peers, and any $m$ of them can collectively authorize key rotation.
This eliminates the single point of failure of traditional key management
while preserving decentralization.


\chapter{Governance}
\label{ch:governance}

\section{Overview}

The Governance cluster implements the decision-making infrastructure. Its 7
zomes (plus 1 bridge) provide a complete governance toolkit:

\begin{enumerate}
  \item \textbf{proposals} --- Proposal creation, amendment, and lifecycle management.
  \item \textbf{voting} --- Multiple voting mechanisms (simple majority, ranked
    choice, quadratic, conviction).
  \item \textbf{threshold-signing} --- Distributed Key Generation (DKG) for
    threshold signatures on governance decisions.
  \item \textbf{councils} --- Council formation, membership management, and
    delegation.
  \item \textbf{constitution} --- Constitutional document management with
    amendment procedures.
  \item \textbf{execution} --- Automated execution of approved governance
    decisions.
  \item \textbf{jurisdiction} --- Jurisdictional boundary management for
    nested governance.
  \item \textbf{bridge} --- Cross-cluster governance dispatch.
\end{enumerate}

\section{Proposal Lifecycle}

\begin{enumerate}
  \item \textbf{Draft}: Any Participant+ can create a proposal draft.
  \item \textbf{Discussion}: The community discusses and suggests amendments.
  \item \textbf{Voting}: Once discussion period ends, voting opens. The voting
    mechanism depends on the proposal type (configured per jurisdiction).
  \item \textbf{Resolution}: Votes are tallied. Trust-weighted vote
    totals determine the outcome.
  \item \textbf{Execution}: Approved proposals trigger automated execution via
    the execution zome, which dispatches the appropriate cross-cluster calls.
\end{enumerate}

\section{Distributed Key Generation}

The threshold-signing zome implements a DKG protocol that produces shared keys
without any single party knowing the complete private key. This enables:

\begin{itemize}
  \item Multi-signature treasury operations (Finance cluster integration)
  \item Constitutional amendment signing (requires Guardian-tier quorum)
  \item Emergency response authorization (distributed circuit-breaker)
\end{itemize}

\section{Constitutional Governance}

The constitution zome manages the foundational rules of a Mycelix community.
Constitutional documents are versioned, and amendments require:

\begin{itemize}
  \item Proposer: Steward tier minimum
  \item Identity dimension: $\geq 0.75$ (HighlyAssured)
  \item Community dimension: $\geq 0.5$
  \item Supermajority vote (configurable, default 67\%)
  \item Mandatory discussion period (configurable, default 14 days)
\end{itemize}


\chapter{Finance}
\label{ch:finance}

\section{Three-Currency Economics}

The Finance cluster implements a three-currency system designed to serve
different economic functions:

\begin{description}
  \item[SAP (Sapience)] A stable unit of account pegged to a basket of real
    goods and services. Used for pricing, contracts, and accounting. Cannot be
    created ex nihilo---minted only through verified economic activity. Subject
    to demurrage, with compost redistributed to commons pools as a maintenance
    flow rather than a hoarding subsidy.
  \item[TEND (Tenderness)] A mutual credit currency for care work, community
    service, and non-market contributions. Zero-sum by construction: every
    positive balance is matched by an equal negative balance elsewhere in the
    network. TEND does not implement demurrage.
  \item[MYCEL (Mycelium)] A commons currency representing collective
    stewardship. Used for treasury operations, public goods funding, and
    inter-community exchange.
\end{description}

\section{Cluster Zomes}

\begin{enumerate}
  \item \textbf{currency-mint} --- Currency creation and supply management.
  \item \textbf{payments} --- Transaction processing for all three currencies.
  \item \textbf{tend} --- TEND-specific mutual credit operations.
  \item \textbf{treasury} --- Trust-gated treasury management.
  \item \textbf{staking} --- MYCEL staking for governance weight.
  \item \textbf{recognition} --- Non-monetary recognition and gratitude tokens.
  \item \textbf{price-oracle} --- Decentralized price discovery for SAP basket.
  \item \textbf{bridge} --- Cross-cluster financial operations.
\end{enumerate}

\section{Trust-Gated Treasury}

Treasury operations are gated by governance tier:

\begin{itemize}
  \item \textbf{View treasury balance}: Observer+
  \item \textbf{Propose expenditure}: Citizen+
  \item \textbf{Approve expenditure}: Steward+ (threshold signature required)
  \item \textbf{Emergency fund release}: Guardian (DKG multi-sig)
\end{itemize}

\section{Demurrage and Composting}

SAP implements demurrage---a periodic reduction in idle balance that reflects
the carrying cost of the real basket backing the currency and discourages
passive hoarding. This is inspired by Silvio Gesell's Freigeld (``free money'')
concept and implemented as:

\begin{equation}
  B(t) = B(t_0) \cdot e^{-\lambda (t - t_0)}
\end{equation}

where $\lambda$ is the demurrage rate. The ``composted'' value (the amount lost
to demurrage) is routed to commons pools, completing the metabolic cycle.

TEND, by contrast, is mutual credit rather than a composting currency. Its
stability comes from bounded credit lines, reciprocity, and zero-sum settlement,
not from timed balance decay.

Currency composting is the explicit mechanism by which retired or excess currency
units are removed from circulation. It is idempotent: composting an already-composted
unit is a no-op, and the total supply is always zero-sum (minted $-$ composted $=$
circulating).

\section{Oracle and Price Discovery}

The price-oracle zome implements decentralized price feeds for the SAP basket.
Multiple oracle agents submit price attestations, and an accuracy-weighted
trimmed median is used as the canonical price. Oracle submissions require
Citizen tier minimum, and when there is a prior consensus the oracle compares
new consensus against the previous stored median to detect volatility and
trigger counter-cyclical TEND limit expansion.
contract includes protections against:

\begin{itemize}
  \item Division by zero (NaN prices)
  \item Stale prices (TTL enforcement)
  \item Manipulation (outlier detection, reputation-weighted aggregation)
\end{itemize}


\chapter{Commons}

\section{The Largest Cluster}

The Commons cluster, with 39 zomes across 9 domains (plus mesh-time and
resource-mesh), is the largest cluster in Mycelix. It implements the practical
infrastructure of daily community life.

\section{Domain Overview}

\subsection{Property (4 zomes)}

\begin{itemize}
  \item \textbf{property\_registry}: Land and building registration with
    provenance tracking.
  \item \textbf{property\_transfer}: Ownership transfer with trust-gated
    approval.
  \item \textbf{property\_disputes}: Structured dispute resolution for property
    conflicts.
  \item \textbf{property\_commons}: Community land trust management.
\end{itemize}

\subsection{Housing (6 zomes)}

Housing cooperatives are first-class citizens in Mycelix:

\begin{itemize}
  \item \textbf{housing\_units}: Unit registry and availability tracking.
  \item \textbf{housing\_membership}: Cooperative membership management.
  \item \textbf{housing\_finances}: Shared housing financial operations.
  \item \textbf{housing\_maintenance}: Maintenance request and scheduling.
  \item \textbf{housing\_clt}: Community Land Trust governance.
  \item \textbf{housing\_governance}: Housing-specific governance (distinct from
    the system-level Governance cluster).
\end{itemize}

\subsection{Care (5 zomes)}

\begin{itemize}
  \item \textbf{care\_timebank}: Time-banking for care work, integrated with
    TEND currency.
  \item \textbf{care\_circles}: Care circle formation and coordination.
  \item \textbf{care\_matching}: Matching care needs with available caregivers.
  \item \textbf{care\_plans}: Structured care planning for individuals and families.
  \item \textbf{care\_credentials}: Caregiver credential verification.
\end{itemize}

\subsection{Mutual Aid (7 zomes)}

\begin{itemize}
  \item \textbf{mutualaid\_needs}: Need posting and categorization.
  \item \textbf{mutualaid\_circles}: Mutual aid circle management.
  \item \textbf{mutualaid\_governance}: Circle-level governance.
  \item \textbf{mutualaid\_pools}: Resource pooling for collective needs.
  \item \textbf{mutualaid\_requests}: Request processing and fulfillment tracking.
  \item \textbf{mutualaid\_resources}: Resource inventory and availability.
  \item \textbf{mutualaid\_timebank}: Mutual aid time banking.
\end{itemize}

\subsection{Water (5 zomes)}

\begin{itemize}
  \item \textbf{water\_flow}: Water flow monitoring and allocation.
  \item \textbf{water\_purity}: Water quality testing and reporting.
  \item \textbf{water\_capture}: Rainwater harvesting coordination.
  \item \textbf{water\_steward}: Watershed stewardship governance.
  \item \textbf{water\_wisdom}: Traditional and scientific water knowledge.
\end{itemize}

\subsection{Food (4 zomes)}

\begin{itemize}
  \item \textbf{food\_production}: Community garden and farm management.
  \item \textbf{food\_distribution}: Food distribution coordination.
  \item \textbf{food\_preservation}: Preservation knowledge and scheduling.
  \item \textbf{food\_knowledge}: Seed libraries, recipes, growing guides.
\end{itemize}

\subsection{Transport (3 zomes)}

\begin{itemize}
  \item \textbf{transport\_routes}: Community transport route planning.
  \item \textbf{transport\_sharing}: Vehicle and ride sharing coordination.
  \item \textbf{transport\_impact}: Environmental impact tracking.
\end{itemize}

\subsection{Cross-Cutting: Mesh-Time and Resource-Mesh}

Two cross-cutting zomes serve the entire Commons cluster:

\begin{itemize}
  \item \textbf{mesh-time}: Coordinated time-keeping across domains, enabling
    temporal queries (``what water quality readings existed when this food was
    distributed?'').
  \item \textbf{resource-mesh}: Cross-domain resource allocation and tracking,
    enabling a community garden to request water allocation from the water
    domain.
\end{itemize}

\section{Testing: 5,276 Tests}

The Commons cluster is tested with 5,276 automated tests---the highest count
of any single cluster. These include unit tests for each zome, integration
tests for cross-zome operations, and sweettest conductor tests for full
Holochain runtime validation.


\chapter{Hearth: Family and Community}

\epigraph{The family is one of nature's masterpieces.}{George Santayana,
\textit{The Life of Reason} (1905)}

\section{Why a Family Cluster?}

Governance frameworks typically model individuals as atomistic rational actors.
But humans live in families, households, and kinship networks. Care work, child
rearing, elder support, and collective decision-making within families are among
the most consequential governance activities in any community---yet they are
invisible to most civic infrastructure.

The Hearth cluster (12 zomes, 1,023 tests) brings family and community
relationships into Mycelix as first-class governance primitives.

\section{Domain Overview}

\subsection{Kinship Networks (2 zomes)}

The kinship zomes model family and chosen-family relationships:

\begin{itemize}
  \item \textbf{kinship\_registry}: Registration of kinship bonds (biological,
    adoptive, chosen-family) with mutual consent. Both parties must confirm
    the relationship before it is committed to the DHT.
  \item \textbf{kinship\_governance}: Family-level governance---shared
    decision-making weighted by kinship role and trust. Parents, elders, and
    designated guardians receive contextual governance weight within the kinship
    network.
\end{itemize}

\subsection{Gratitude Exchange (1 zome)}

The gratitude zome implements a non-monetary recognition system for care work,
emotional support, and community contribution. Gratitude tokens:

\begin{itemize}
  \item Are non-transferable (cannot be hoarded or traded)
  \item Feed the community dimension of the sender's and receiver's trust profiles
  \item Decay with time (30-day half-life), preventing gratitude inflation
  \item Are visible to the kinship network but private to outsiders
\end{itemize}

\subsection{Care Coordination (1 zome)}

The care zome coordinates within-family care responsibilities:

\begin{itemize}
  \item Scheduling shared care duties (childcare, elder care, disability support)
  \item Matching care needs with available family members
  \item Integration with the Commons cluster's care timebank for extended
    community care
\end{itemize}

\subsection{Autonomy Tracking (1 zome)}

A distinctive feature of Hearth: the autonomy zome tracks the developing
autonomy of minors and dependents. As a child grows, their governance
capability within the family (and eventually the broader community) increases
progressively. This mirrors the trust tier system at the family scale:

\begin{itemize}
  \item Young children: Observer (family decisions are visible to them)
  \item Older children: Participant (can express preferences)
  \item Adolescents: Citizen (vote weight in family governance)
  \item Adults: Full autonomy (Steward/Guardian eligible)
\end{itemize}

\subsection{Collective Decisions (1 zome)}

Family-level collective decisions use trust-weighted voting where weights
are determined by kinship role, age, and demonstrated engagement with the
family's governance. A family deciding where to live, how to allocate a
shared resource, or how to respond to a community need uses the same
governance primitives as the system-level Governance cluster.

\subsection{Story Archiving (1 zome)}

The stories zome preserves family narratives, oral histories, and cultural
knowledge. Stories are:

\begin{itemize}
  \item Stored on the author's source chain (tamper-evident)
  \item Linked to kinship relationships (``grandmother's recipe'')
  \item Optionally shared with the broader community or the Knowledge cluster
  \item Tagged with epistemic classification (E0-N2-M2 for cultural knowledge)
\end{itemize}

\subsection{Milestones, Rhythms, Emergency, and Resources (4 zomes)}

\begin{description}
  \item[Milestones] Track family events (births, graduations, ceremonies) as
    governance-relevant markers that can trigger tier transitions.
  \item[Rhythms] Model recurring family patterns (meal schedules, care
    rotations) that the mesh-time system can synchronize across households.
  \item[Emergency] Family-level emergency coordination, integrated with the
    Civic cluster's emergency system for escalation.
  \item[Resources] Shared family resource tracking (tools, vehicles, spaces),
    with lending and return workflows.
\end{description}

\subsection{Bridge (1 zome)}

The Hearth bridge dispatches cross-cluster requests---most commonly to
Identity (for kinship credential verification), Commons (for care coordination),
and Governance (for family-level proposal escalation to community governance).

\section{Trust-Weighted Family Governance}

Family governance in Hearth is trust-weighted but with different dimension
priorities than system-level governance. The community dimension (peer trust)
carries the highest weight within a kinship network, because family trust is
fundamentally relational rather than transactional.

The 1,023 tests cover kinship lifecycle (creation, modification, dissolution),
gratitude exchange mechanics, autonomy progression, and cross-cluster bridge
dispatch. 48 sweettest integration tests verify full Holochain runtime behavior,
including 9 trust-gating tests that validate tier requirements for
family governance actions.


\chapter{Personal and Attribution}

\section{Personal: Private Vaults}

The Personal cluster (4 zomes including bridge, 20 tests) implements
private data management:

\subsection{Identity Vault}

The identity vault stores the agent's most sensitive identity artifacts:

\begin{itemize}
  \item Recovery seed phrases (encrypted, never DHT-replicated)
  \item Biometric templates (stored locally, used for MFA verification)
  \item Delegation keys (sub-passports for authorized devices)
  \item Identity dimension history (track MFA level changes over time)
\end{itemize}

The vault is strictly source-chain-only: entries are never committed to the
DHT, ensuring that even a compromised DHT neighborhood cannot access the
agent's core identity material.

\subsection{Health Vault}

The health vault stores personal health data under patient control:

\begin{itemize}
  \item Medical records (FHIR-formatted, from the Health cluster)
  \item Prescription history
  \item Genetic data (if voluntarily submitted)
  \item Mental health assessments
\end{itemize}

Data sharing is opt-in and granular: an agent can share lab results with
their care circle (via the Hearth cluster) without exposing their full
medical history.

\subsection{Credential Wallet}

The credential wallet manages the agent's verifiable credentials:

\begin{itemize}
  \item W3C Verifiable Credentials from the Identity cluster
  \item Educational credentials from the EduNet cluster
  \item Care credentials from the Commons cluster
  \item Governance tier credentials (ConsciousnessCredentials)
\end{itemize}

The wallet presents credentials selectively---an agent proving they are
Citizen tier need not reveal their identity dimension score, only the derived
tier.

\section{Attribution: Tracking Reciprocity}

The Attribution cluster (3 zomes, 17 tests) provides cross-cutting
infrastructure for tracking who contributed what to whom:

\begin{description}
  \item[Dependency Registry] Tracks dependency relationships between
    creative works, software, datasets, and community resources.
    When an educational curriculum uses a Knowledge cluster claim,
    the dependency is recorded.
  \item[Usage Receipts] Immutable records of resource usage. When the
    Music cluster plays a track, a usage receipt is created on the
    listener's source chain and replicated to the artist's DHT neighborhood.
  \item[Reciprocity Engine] Computes net reciprocity scores between
    agents and communities. An agent who consumes many community
    resources but contributes little will see their engagement dimension
    decrease. An agent who contributes generously will see it increase.
\end{description}

Attribution integrates with Finance (TEND payments for usage), Music (artist
royalties), Education (curriculum attribution), and Knowledge (citation
tracking).


%% ════════════════════════════════════════════════════════════════════════════
%% PART IV: DOMAIN CLUSTERS
%% ════════════════════════════════════════════════════════════════════════════

\part{Domain Clusters}

\chapter{Health}

\section{15-Zome Architecture}

The Health cluster is Mycelix's most medically comprehensive component, with 15
active zomes organized in two tiers:

\textbf{Tier 1 (MVP):}
\begin{enumerate}
  \item Patient records management
  \item Provider registry and credentialing
  \item Appointment scheduling
  \item Prescription management
  \item Lab results tracking
  \item Referral coordination
  \item Billing and insurance claims
\end{enumerate}

\textbf{Tier 2 (Advanced):}
\begin{enumerate}
  \item \textbf{FHIR} --- HL7 FHIR resource modeling for healthcare interoperability.
  \item \textbf{CDS} --- Clinical Decision Support with tier-gated
    recommendation generation.
  \item \textbf{Telehealth} --- Decentralized telemedicine session management.
  \item \textbf{Trials} --- Clinical trial coordination and data collection.
  \item \textbf{Insurance} --- Decentralized mutual health insurance pools.
  \item \textbf{Nutrition} --- Nutritional planning and dietary management.
  \item \textbf{Mental health} --- Mental health assessment and tracking.
  \item \textbf{Public health} --- Epidemiological surveillance and reporting.
\end{enumerate}

\section{FHIR Interoperability}

The FHIR (Fast Healthcare Interoperability Resources) zome implements HL7 FHIR
R4 resource types as Holochain entries. This enables:

\begin{itemize}
  \item Data exchange with existing healthcare systems
  \item Standardized patient records across Mycelix communities
  \item Regulatory compliance with healthcare data standards
\end{itemize}

\section{Privacy Architecture}

Health data requires the strongest privacy protections. The Health cluster
implements:

\begin{itemize}
  \item Patient-controlled data sharing (source chain privacy)
  \item Encrypted DHT entries for sensitive records
  \item Tier-gated access (provider credentialing required)
  \item Audit trails on all data access
\end{itemize}

\section{Clinical Decision Support}

The CDS (Clinical Decision Support) zome is tier-gated at two levels:

\begin{enumerate}
  \item \textbf{Provider access}: Only agents with verified healthcare
    credentials (from the Identity cluster's education and verifiable
    credential zomes) can query the CDS system.
  \item \textbf{Recommendation generation}: CDS recommendations are tagged
    with the governance tier of the requesting provider and the evidence
    level of the recommendation. A Participant-tier provider receives only
    well-established guidelines, while a Steward-tier provider can access
    experimental protocols and off-label usage data.
\end{enumerate}

This tiered access prevents information overload for less-experienced
providers while ensuring that experienced practitioners have access to
the full knowledge base.

\section{Mutual Health Insurance}

The Insurance zome implements decentralized mutual health insurance pools:

\begin{itemize}
  \item Community members contribute to a shared pool (denominated in TEND
    or MYCEL)
  \item Claims are submitted with medical evidence (from the patient records
    zome) and adjudicated by a panel of Citizen+ tier pool members
  \item Large claims require Steward-level approval and DKG multi-signature
    from the Governance cluster
  \item Pool governance follows the same trust-weighted model as
    system-level governance
\end{itemize}

This model eliminates the profit motive of commercial insurance while
maintaining the risk-pooling benefits.


\chapter{Knowledge}
\label{ch:knowledge}

\section{Epistemic Infrastructure}

The Knowledge cluster implements decentralized knowledge verification---an
``epistemic mesh'' where claims are verified through multiple independent
pathways:

\begin{enumerate}
  \item \textbf{claims} --- Structured claim registration with provenance metadata.
  \item \textbf{graph} --- Knowledge graph construction with typed relationships.
  \item \textbf{query} --- Semantic query processing across the knowledge graph.
  \item \textbf{inference} --- Logical inference engine for deriving new claims
    from existing ones.
  \item \textbf{factcheck} --- Multi-source fact verification with confidence scoring.
  \item \textbf{markets\_integration} --- Prediction market integration for
    epistemic value discovery.
  \item \textbf{bridge} --- Cross-cluster knowledge dispatch.
\end{enumerate}

\section{Claim Lifecycle}

\begin{enumerate}
  \item \textbf{Registration}: An agent registers a claim with supporting evidence.
    The claim entry includes: statement, evidence links, confidence, domain tags.
  \item \textbf{Verification}: Independent agents verify or dispute the claim.
    Verification quality is weighted by the verifier's governance tier and
    domain engagement.
  \item \textbf{Inference}: The inference engine derives implications from
    verified claims, extending the knowledge graph automatically.
  \item \textbf{Factchecking}: The factcheck zome aggregates verification
    results into a confidence score.
  \item \textbf{Market pricing}: Prediction markets assign economic value to
    claim accuracy, creating incentives for truthful reporting.
\end{enumerate}

\section{Knowledge Graph}

The graph zome maintains a typed knowledge graph where nodes are claims and
edges are relationships (supports, contradicts, extends, specializes). The
graph is distributed across the DHT, with each neighborhood validating the
relationships in its shard.

Query processing supports:
\begin{itemize}
  \item Path queries (``what claims support this conclusion?'')
  \item Subgraph extraction (``all claims in the climate domain'')
  \item Confidence propagation (``what is the transitive confidence in this claim?'')
\end{itemize}


\section{Prediction Markets}

The markets\_integration zome connects knowledge claims to economic incentives:

\begin{enumerate}
  \item A Citizen+ agent creates a prediction market for a specific, verifiable
    claim with a defined resolution date.
  \item Participants stake MYCEL on the outcome (true/false/probability).
  \item Market prices reflect the collective's confidence in the claim.
  \item Upon resolution, correct predictors are rewarded proportionally,
    and their reputation dimension increases.
  \item Incorrect predictions reduce reputation, creating accountability.
\end{enumerate}

The prediction market mechanism provides a ``skin in the game'' complement to
the expert verification pathway. Claims that have both high expert verification
scores and high market confidence are the most robustly validated.

\section{Integration Patterns}

The Knowledge cluster integrates with other clusters through several patterns:

\begin{description}
  \item[Health $\to$ Knowledge] Medical claims from clinical trials feed
    the knowledge graph, enabling meta-analysis across decentralized
    research sites.
  \item[Climate $\to$ Knowledge] Carbon credit verification claims are
    registered in the knowledge graph and can be independently verified
    through satellite observation data from the Space cluster.
  \item[Supply Chain $\to$ Knowledge] Provenance claims from the supply
    chain are cross-referenced with production data from the Commons
    cluster's food and water domains.
  \item[Education $\to$ Knowledge] Educational content references knowledge
    graph claims, ensuring that curricula are backed by verified knowledge.
\end{description}


\chapter{Space and Climate}

\section{Space: Orbital Stewardship}

The Space cluster addresses a growing governance challenge: who manages orbital
space? With 5 zomes plus an orbital mechanics library, it implements:

\begin{enumerate}
  \item \textbf{orbital\_objects} --- Satellite and debris tracking registry.
  \item \textbf{observations} --- Decentralized observation reporting from
    ground stations and amateur astronomers.
  \item \textbf{conjunctions} --- Collision probability computation and warning
    dissemination.
  \item \textbf{debris\_bounties} --- Economic incentives for debris removal,
    funded through MYCEL commons currency.
  \item \textbf{traffic\_control} --- Decentralized orbital traffic management
    protocols.
\end{enumerate}

The orbital mechanics library provides Keplerian propagation, conjunction
assessment, and maneuver planning---the computational backbone for governance
decisions about orbital resource allocation.

\section{Climate: Carbon and Beyond}

The Climate cluster manages environmental stewardship:

\begin{enumerate}
  \item \textbf{carbon} --- Carbon credit registry with tier-gated
    verification. Credits are issued for verified carbon sequestration and
    can be traded or retired.
  \item \textbf{projects} --- Climate project registration, monitoring, and
    impact reporting.
  \item \textbf{bridge} --- Cross-cluster climate data dispatch (e.g., carbon
    impact of supply chain operations).
\end{enumerate}

\section{Integration: Space-Climate Nexus}

Space and climate governance intersect in satellite-based Earth observation:
the Space cluster's observation network can verify climate project claims
(deforestation monitoring, agricultural carbon sequestration) through the
Knowledge cluster's factcheck pipeline.

\section{Debris Economics}

The debris\_bounties zome creates economic incentives for orbital cleanup:

\begin{enumerate}
  \item Debris objects are cataloged in the orbital\_objects registry with
    estimated removal difficulty and collision risk.
  \item The community allocates MYCEL bounties proportional to the object's
    collision risk (computed by the conjunctions zome).
  \item Removal operators claim bounties by submitting observation evidence
    that the debris has been deorbited or moved to a graveyard orbit.
  \item Claims are verified through the Knowledge cluster's factchecking
    pipeline using independent observations from multiple ground stations.
\end{enumerate}

This mechanism applies the same trust-weighted governance to space
stewardship: bounty creation requires Citizen+ tier, bounty approval
requires Steward+, and emergency collision avoidance can be triggered
by Guardian-tier agents.

\section{Carbon Credit Lifecycle}

The carbon zome implements a full carbon credit lifecycle:

\begin{enumerate}
  \item \textbf{Project registration}: A climate project (reforestation,
    renewable energy, soil carbon) is registered with baseline measurements.
  \item \textbf{Monitoring}: Periodic measurement reports are submitted
    with evidence links (satellite data, sensor readings, third-party audits).
  \item \textbf{Credit issuance}: Verified sequestration generates carbon
    credits, recorded as Holochain entries on the project operator's source
    chain.
  \item \textbf{Trading}: Credits can be transferred between agents, retired
    (permanently removed from circulation), or used to offset emissions.
  \item \textbf{Retirement}: Retired credits are composted (similar to
    currency composting in the Finance cluster), ensuring they cannot be
    double-counted.
\end{enumerate}

All operations are tier-gated: credit issuance requires Citizen+
with verified engagement in the climate domain, and credit retirement is
auditable through the source chain.


\chapter{Music and Education}

\section{Music: Attribution-First Royalties}

The Music cluster (4 zomes plus 14 support crates) implements a
royalty-free-by-default, attribution-required music platform:

\begin{enumerate}
  \item \textbf{catalog} --- Music catalog management with provenance tracking.
  \item \textbf{plays} --- Play count tracking and listener analytics.
  \item \textbf{balances} --- Artist balance management (TEND-denominated).
  \item \textbf{trust} --- Artist verification and anti-plagiarism.
\end{enumerate}

The Attribution cluster (\texttt{mycelix-attribution}) provides the cross-cutting
infrastructure: dependency registry, usage receipts, and reciprocity tracking.

\section{Education: Curriculum and Credentials}

The Education cluster (10 zomes, \texttt{mycelix-edunet}) implements:

\begin{itemize}
  \item Curriculum design and sharing
  \item Learner progress tracking
  \item Peer assessment with tier-gated credentialing
  \item Educational credential issuance (integrated with Identity cluster)
  \item Knowledge graph integration for learning path recommendation
\end{itemize}


\chapter{Civic Infrastructure}
\label{ch:civic}

\epigraph{Justice is what love looks like in public.}{Cornel West}

\section{17 Zomes for Civic Life}

The Civic cluster (18 zomes including bridge, 2,273 tests) implements three
domains of civic infrastructure: restorative justice, emergency coordination,
and decentralized media. Together they address the functions that nation-states
perform through courts, emergency services, and regulated press---reimagined
for a decentralized, trust-weighted governance model.

\section{Restorative Justice}

The Justice domain explicitly rejects the punitive model of criminal records,
incarceration, and permanent stigma. Instead, it implements restorative
justice---a process-oriented approach that seeks to repair harm rather than
punish offenders.

\subsection{Process, Not Punishment}

The justice zomes implement a structured conflict resolution lifecycle:

\begin{enumerate}
  \item \textbf{Incident report}: Any Participant+ agent can file a structured
    incident report describing harm experienced. Reports are stored on the
    reporter's source chain and shared with a mediator pool.
  \item \textbf{Mediator assignment}: Trained mediators (Citizen+ tier with
    justice domain engagement) are assigned based on availability, domain
    expertise, and absence of conflict of interest.
  \item \textbf{Circle process}: The mediator convenes a restorative circle
    including the affected parties, community supporters, and (optionally)
    the person who caused harm. The process follows structured dialogue
    protocols.
  \item \textbf{Agreement}: The circle produces a restorative agreement---a
    set of commitments by the responsible party (restitution, community
    service, behavioral change) with a timeline and verification criteria.
  \item \textbf{Monitoring}: Agreement compliance is tracked through the
    governance system. Completion feeds the reputation dimension positively;
    non-compliance triggers graduated sanctions (reputation decay, tier
    reduction).
\end{enumerate}

\subsection{No Permanent Records}

A critical design decision: completed restorative processes are
\emph{archived} rather than permanently attached to an agent's identity.
After a configurable period (default: 1 year post-completion), the
incident hash is removed from the agent's public profile links. The
source chain entry remains (Holochain source chains are append-only),
but it is no longer discoverable through DHT queries. This implements
the ``right to rehabilitation'' in code.

\section{Emergency Coordination}

The Emergency domain provides infrastructure for disaster response,
community crises, and urgent coordination:

\begin{description}
  \item[Incidents] Structured incident reporting with severity levels,
    geolocation, and resource requirements. High-severity incidents
    trigger automatic cross-cluster notifications.
  \item[Triage] Resource triage using a priority queue weighted by
    severity, affected population, and available resources. Triage
    decisions require Citizen+ tier and are auditable.
  \item[Shelters] Shelter registration, capacity tracking, and
    occupancy management. Integrated with the Commons cluster's housing
    domain for emergency housing allocation.
  \item[Communications] Emergency broadcast system using the bridge
    dispatch to reach all clusters simultaneously. Guardian-tier agents
    can trigger network-wide emergency alerts. When the mesh radio
    bridge is available (feature \texttt{mesh}), emergency alerts
    propagate via LoRa radio even when internet connectivity is lost.
\end{description}

\section{Decentralized Media}

The Media domain implements a decentralized publishing and curation platform:

\begin{description}
  \item[Publication] Content publishing with provenance tracking. Every
    published piece is signed and timestamped on the author's source
    chain, providing tamper-evident authorship.
  \item[Curation] Community-driven content curation using trust-weighted
    recommendations. Higher-tier curators' selections carry more weight,
    but all Participant+ agents can curate.
  \item[Fact-checking] Integration with the Knowledge cluster's factcheck
    pipeline. Media claims can be tagged for verification, and the
    verification status is displayed alongside the content.
  \item[Attribution] Integration with the Attribution cluster for proper
    crediting of sources, quotes, and derivative works.
  \item[Resonance Feed] A trust-weighted content feed that surfaces
    content based on community trust signals rather than engagement
    metrics. This explicitly rejects the attention-economy model that
    rewards outrage and sensationalism.
\end{description}

\section{Testing: 2,273 Tests}

The Civic cluster's 2,273 tests cover justice process lifecycle, emergency
coordination workflows, media publication and curation, and cross-cluster
bridge dispatch. 17 sweettest integration tests verify full conductor-level
behavior, including cross-domain interactions (e.g., an emergency incident
triggering a justice mediation process).


\chapter{Mail, Supply Chain, and Marketplace}

\section{Mail: PQC-Encrypted Decentralized Email}

The Mail cluster (12 zomes) implements a fully decentralized email system with
post-quantum cryptographic security:

\begin{description}
  \item[Key exchange] ML-KEM-768 (NIST FIPS 203) key encapsulation for
    establishing shared secrets between correspondents.
  \item[Message encryption] AES-256-GCM symmetric encryption using the
    ML-KEM-derived shared secret.
  \item[Mailbox management] Each agent maintains a private mailbox on their
    source chain. Incoming messages are encrypted to the recipient's public
    key and stored in their DHT neighborhood.
  \item[Address book] Decentralized address resolution using the Identity
    cluster's name-registry zome.
  \item[Attachments] Large attachments are chunked, encrypted, and distributed
    across the DHT with content-addressed hashes.
  \item[Spam prevention] Trust-weighted message filtering: messages from
    agents below Observer tier are quarantined by default. Community
    trust dimension modulates filter sensitivity.
\end{description}

The 12 zomes cover the full email lifecycle: compose, encrypt, send, receive,
decrypt, reply, forward, search, folder management, contact management,
attachment handling, and bridge dispatch. All message content is
end-to-end encrypted; even DHT validators cannot read message bodies.

\section{Supply Chain: Provenance Tracking}

The Supply Chain cluster (8 zomes) implements end-to-end provenance tracking
for physical goods:

\begin{description}
  \item[Inventory] Product registration with unique identifiers, batch
    tracking, and location history.
  \item[Logistics] Shipment management, route optimization, and delivery
    confirmation with cryptographic proof-of-delivery.
  \item[Procurement] Purchase order management with trust-gated supplier
    approval (Citizen+ tier required for high-value procurement).
  \item[Verification] Product authenticity verification using source chain
    provenance---each handler's signature provides a link in the
    chain of custody.
  \item[Trust] Supplier reputation tracking integrated with the system-wide
    reputation dimension.
  \item[Claims] Dispute resolution for supply chain disagreements, integrated
    with the Civic cluster's justice domain.
  \item[Payments] Supply chain payment processing using SAP (for contracted
    prices) and TEND (for informal exchanges).
  \item[Bridge] Cross-cluster dispatch for supply chain queries from other
    clusters (e.g., Food domain querying provenance of ingredients).
\end{description}

\section{Marketplace: Decentralized Commerce with Arbitration}

The Marketplace cluster (8 zomes) provides a decentralized commerce platform:

\begin{itemize}
  \item Product and service listing with tier-gated seller verification
  \item Escrow-based transactions using the Finance cluster's payment
    infrastructure
  \item Buyer and seller reputation tracking
  \item Structured arbitration process for disputes, escalating to the
    Civic cluster's justice domain when marketplace-level resolution fails
  \item Trust-weighted product reviews (higher-tier reviewers' assessments
    carry more weight)
  \item Integration with the Supply Chain cluster for physical goods
    provenance verification
\end{itemize}


\chapter{Energy and DeSci}

\section{Energy: Community Power}

The Energy cluster (5 zomes) manages community-scale energy production,
investment, and distribution:

\begin{description}
  \item[Projects] Energy project registration---solar installations, wind
    turbines, micro-hydro, community batteries---with capacity,
    generation forecasts, and maintenance schedules.
  \item[Investments] Trust-gated community investment in energy projects.
    Investment proposals require Citizen+ tier. Large investments
    (above a configurable threshold) require Steward+ approval and
    DKG multi-signature from the Governance cluster.
  \item[Regenerative] Regenerative energy practices tracking---soil
    carbon sequestration from biochar, forest management, and agricultural
    energy recovery.
  \item[Grid] Community microgrid management: energy balance tracking,
    peer-to-peer energy trading (denominated in MYCEL), and demand
    response coordination.
  \item[Bridge] Cross-cluster dispatch for energy data queries from
    Climate (carbon offsets from renewable generation), Finance
    (energy investment returns), and Commons (resource-mesh integration).
\end{description}

\section{DeSci: Decentralized Science}

The DeSci cluster (\texttt{mycelix-desci}) takes a different architectural
approach from other Mycelix clusters: it implements a REST API (using
Actix-web) rather than Holochain zomes. This design choice reflects the
DeSci cluster's role as a bridge between the Mycelix governance network
and the existing scientific research ecosystem.

\subsection{REST API Architecture}

The DeSci REST API provides:

\begin{itemize}
  \item Research project registration and metadata management
  \item Dataset registration with provenance tracking
  \item Peer review coordination with trust-weighted reviewer assignment
  \item Reproducibility tracking (linking research claims to data,
    code, and computational environment)
  \item Integration with the Knowledge cluster for claim verification
  \item Open access publishing with attribution tracking
\end{itemize}

\subsection{141 Integration Tests}

The DeSci cluster is validated by 141 integration tests that exercise
the full REST API surface, including:

\begin{itemize}
  \item Project lifecycle (create, update, archive)
  \item Dataset upload, versioning, and access control
  \item Review assignment, submission, and aggregation
  \item Cross-cluster queries (Knowledge graph, Attribution)
  \item Authentication and authorization (trust profile verification)
\end{itemize}

\subsection{Federated Learning Integration}

DeSci integrates with the \texttt{mycelix-core} federated learning framework
for privacy-preserving collaborative research. Research institutions can
participate in multi-site studies without sharing raw patient data---only
trust-weighted gradient updates are exchanged through the 0TML pipeline.


%% ────────────────────────────────────────────────────────────────────────────
%% Additional Domain Cluster Chapters
%% ────────────────────────────────────────────────────────────────────────────

\chapter{PQC-Encrypted Decentralized Mail}
\label{ch:mail}

\epigraph{The right of the people to be secure in their persons, houses, papers,
and effects, against unreasonable searches and seizures, shall not be
violated.}{Fourth Amendment, U.S.\ Constitution}

\section{Why Decentralized Mail Requires Post-Quantum Cryptography}

Centralized email is the single largest vector for state surveillance and
corporate data harvesting.  The \texttt{mycelix-mail} cluster (12 zomes,
approximately 100,000 lines of Rust) replaces the entire SMTP/IMAP stack
with a Holochain-native messaging system where \emph{no central mail server
exists} and every message is encrypted with lattice-based post-quantum
primitives.

The cryptographic foundation is ML-KEM-768 (NIST FIPS 203), a module-lattice
key encapsulation mechanism that remains secure against known quantum
algorithms.  Forward secrecy is achieved without a dedicated key server:
each correspondence pair performs a double-ratchet key exchange mediated
entirely through DHT entries on the participants' source chains.

\section{Zome Architecture}

The 12 zomes partition the full email lifecycle into orthogonal concerns:

\begin{longtable}{lp{8cm}}
\toprule
\textbf{Zome} & \textbf{Responsibility} \\
\midrule
\endhead
\texttt{compose}       & Message drafting, template management \\
\texttt{encrypt}       & ML-KEM encapsulation, AES-256-GCM bulk encryption \\
\texttt{send}          & DHT neighborhood publication, delivery receipts \\
\texttt{receive}       & Incoming message polling, decryption delegation \\
\texttt{decrypt}       & Private-key operations via LAIR keystore \\
\texttt{threads}       & Conversation threading, in-reply-to resolution \\
\texttt{attachments}   & Content-addressed chunking, encrypted distribution \\
\texttt{contacts}      & Address book, Identity cluster name-registry integration \\
\texttt{folders}       & Inbox/sent/archive/custom folder management \\
\texttt{search}        & Encrypted search over ciphertext (trapdoor-based) \\
\texttt{spam}          & Trust-weighted spam filtering (see Section~\ref{sec:mail-spam}) \\
\texttt{bridge}        & Cross-cluster dispatch for mail queries \\
\bottomrule
\end{longtable}

\section{Forward Secrecy Without a Key Server}

Classical forward secrecy relies on ephemeral Diffie--Hellman exchanges
coordinated by a trusted server.  Mycelix eliminates the server by storing
ephemeral public keys as \emph{private DHT entries} in the recipient's
neighborhood.  The sender retrieves the latest ephemeral key, encapsulates
a shared secret via ML-KEM, and discards the ephemeral pair after use.  The
result is per-message forward secrecy: compromise of a long-term private key
does not retroactively decrypt past messages.

\section{Encrypted Search Over Ciphertext}

Full-text search is critical for usability but appears incompatible with
end-to-end encryption.  Mycelix implements \emph{trapdoor-based search}:
the sender generates encrypted keyword indices (trapdoors) alongside the
ciphertext.  The recipient's search zome matches query trapdoors against
stored trapdoors without decrypting message bodies.  This provides
substring-level search with zero plaintext exposure to DHT validators.

\section{Spam Filtering and Trust Integration}
\label{sec:mail-spam}

Spam filtering leverages the 8D sovereign profile (Chapter~\ref{ch:trust-profile}):

\begin{itemize}
  \item Messages from agents below Observer tier are quarantined by default.
  \item Community trust dimension modulates filter sensitivity---agents
    vouched for by the recipient's web-of-trust bypass quarantine.
  \item Reputation dimension provides a decay-based spam cost: repeated
    spam reports slash sender reputation ($0.5\times$), eventually triggering
    blacklist ($< 0.05$).
  \item Attachment size limits scale with sender tier (Observer: 1\,MB,
    Guardian: 100\,MB).
\end{itemize}

\section{Thread Management and Attachments}

Conversations are modeled as directed acyclic graphs of DHT entries.
The \texttt{threads} zome maintains link tags for \texttt{in-reply-to},
\texttt{references}, and \texttt{forward-of} relationships.  Large
attachments are split into 256\,KB chunks, individually encrypted, and
distributed across the sender's DHT neighborhood with content-addressed
BLAKE3 hashes.  The recipient reassembles chunks using the manifest
entry attached to the parent message.


\chapter{Supply Chain Provenance}
\label{ch:supplychain}

\epigraph{Trust, but verify.}{Russian proverb, popularized by Ronald Reagan}

\section{End-to-End Provenance Without Consensus Bottleneck}

Traditional supply chain blockchains suffer from the consensus
bottleneck: every transaction must be globally ordered by miners or
validators, limiting throughput and creating single points of failure.
The \texttt{mycelix-supplychain} cluster (8 zomes) uses Holochain's
agent-centric DHT to record provenance \emph{locally} on each handler's
source chain while the DHT's validation rules ensure global consistency
without global consensus.

Each handler in the chain of custody---farmer, processor, shipper,
retailer---signs a \texttt{ProvenanceEntry} on their source chain.
Downstream handlers verify the upstream signature before accepting
custody.  The result is a cryptographically linked chain of custody
that is tamper-evident, auditable, and operates at the throughput of
individual agents rather than global block intervals.

\section{Zome Architecture}

\begin{longtable}{lp{8cm}}
\toprule
\textbf{Zome} & \textbf{Responsibility} \\
\midrule
\endhead
\texttt{inventory}     & Product registration, batch tracking, location history \\
\texttt{logistics}     & Shipment management, route optimization, delivery confirmation \\
\texttt{procurement}   & Purchase orders, tier-gated supplier approval (Citizen+) \\
\texttt{verification}  & Authenticity verification via source chain provenance \\
\texttt{trust}         & Supplier reputation, integrated with system-wide reputation \\
\texttt{claims}        & Dispute resolution, Civic cluster justice integration \\
\texttt{payments}      & SAP (contracted) and TEND (informal) payment processing \\
\texttt{bridge}        & Cross-cluster dispatch (e.g., Food domain querying provenance) \\
\bottomrule
\end{longtable}

\section{Trust-Weighted Supplier Verification}

Supplier verification uses the 8D sovereign profile
(Chapter~\ref{ch:trust-profile}) with domain-specific extensions:

\begin{enumerate}
  \item \textbf{Identity}: Verified business registration via the
    Identity cluster's credential wallet.
  \item \textbf{Reputation}: Delivery reliability score (on-time rate,
    defect rate, dispute history) with temporal decay ($0.998^{\text{days}}$).
  \item \textbf{Community}: Peer endorsements from existing supply chain
    participants, weighted by endorser tier.
  \item \textbf{Engagement}: Volume and frequency of successful
    transactions in the past 90 days.
\end{enumerate}

High-value procurement (above a configurable threshold) requires
Citizen-tier or above for both buyer and supplier, enforced by the
consciousness gating layer in \texttt{bridge-common}.

\section{IoT Sensor Integration for Cold Chain Monitoring}

The logistics zome supports \texttt{SensorReading} entries for
temperature, humidity, GPS coordinates, and shock/vibration events.
IoT gateways publish signed sensor readings to the handler's source
chain at configurable intervals.  Validation rules enforce:

\begin{itemize}
  \item Temperature readings within product-specific acceptable ranges.
  \item GPS coordinates consistent with the declared shipping route.
  \item Timestamp monotonicity (no backdated readings).
  \item Gateway device attestation via the Identity cluster.
\end{itemize}

Cold chain excursions generate \texttt{AlertEntry} records that propagate
to all downstream handlers via DHT links, enabling automated rejection
of compromised shipments before they reach consumers.

\section{Cross-Cluster Integration}

Supply chain provenance integrates with multiple Mycelix clusters:

\begin{longtable}{lp{8cm}}
\toprule
\textbf{Cluster} & \textbf{Integration} \\
\midrule
\endhead
Commons (Food)   & Ingredient provenance queries via bridge dispatch \\
Finance          & SAP/TEND payment settlement for purchase orders \\
Civic (Justice)  & Dispute escalation when claims zome resolution fails \\
Marketplace      & Product authenticity verification for listed goods \\
Identity         & Supplier business credential verification \\
\bottomrule
\end{longtable}


\chapter{The Marketplace and Arbitration}
\label{ch:marketplace}

\epigraph{Markets are conversations.}{The Cluetrain Manifesto}

\section{Decentralized Commerce Without Central Authority}

The \texttt{mycelix-marketplace} cluster (8 zomes) implements a
decentralized marketplace where buyers and sellers transact without
a platform intermediary extracting fees or controlling access.  Seller
verification is trust-weighted: listing a product requires Participant
tier or above, while high-value listings (above a configurable
threshold) require Citizen tier.

\section{Trust-Weighted Seller Ratings}

Seller ratings depart from the centralized five-star model.  Each
review is weighted by the reviewer's combined trust score
(Chapter~\ref{ch:trust-profile}):

\begin{equation}
  R_{\text{weighted}} = \frac{\sum_{i} w_i \cdot r_i}{\sum_{i} w_i},
  \quad w_i = \text{combined\_score}(\text{reviewer}_i)
\end{equation}

This means a Guardian-tier reviewer's assessment carries substantially
more weight than an Observer's.  The temporal decay of the reputation
dimension ($0.998^{\text{days}}$) ensures that sellers must maintain
quality over time---a burst of positive reviews from years ago fades
naturally.

\section{Escrow Without Central Authority}

Marketplace escrow uses the Finance cluster's payment infrastructure
(Chapter~\ref{ch:finance}) with a three-phase commit:

\begin{longtable}{clp{7cm}}
\toprule
\textbf{Phase} & \textbf{Actor} & \textbf{Action} \\
\midrule
\endhead
1 & Buyer   & Funds locked in escrow via SAP payment \\
2 & Seller  & Ships goods; logistics zome records dispatch \\
3a & Buyer  & Confirms receipt $\rightarrow$ escrow releases to seller \\
3b & Timeout & Auto-release after configurable period (default 14 days) \\
3c & Dispute & Either party initiates arbitration (see Section~\ref{sec:arbitration}) \\
\bottomrule
\end{longtable}

The escrow is implemented as a \texttt{LockedPaymentEntry} on the
Finance cluster, referenced by action hash from the Marketplace
listing.  No central escrow service holds funds; the payment is
cryptographically locked until one of the three resolution paths
executes.

\section{Arbitration System for Dispute Resolution}
\label{sec:arbitration}

When buyer and seller cannot resolve a dispute, the Marketplace
cluster's arbitration zome activates a structured resolution process:

\begin{enumerate}
  \item \textbf{Filing}: Either party files a \texttt{DisputeEntry} with
    evidence (transaction records, sensor data, communications).
  \item \textbf{Arbitrator selection}: Three arbitrators are selected from
    Steward+ tier agents with domain engagement in the relevant product
    category.  Selection uses deterministic randomness seeded by the
    dispute entry's action hash to prevent manipulation.
  \item \textbf{Evidence review}: Arbitrators review submitted evidence
    within a configurable window (default 7 days).
  \item \textbf{Ruling}: Majority vote determines the outcome---full
    refund, partial refund, or seller payment.  The ruling is recorded
    on-chain and the escrow is disbursed accordingly.
  \item \textbf{Escalation}: If either party rejects the ruling, the
    dispute escalates to the Civic cluster's justice domain
    (Chapter~\ref{ch:civic}) for community-level adjudication.
\end{enumerate}

\section{Product Categories and Search}

Product listings support hierarchical categories, free-text search,
and tag-based filtering.  Search is implemented via DHT link tags
rather than a centralized search index:

\begin{itemize}
  \item Category links: \texttt{ProductEntry} $\xrightarrow{\text{category}}$ \texttt{CategoryAnchor}
  \item Tag links: \texttt{ProductEntry} $\xrightarrow{\text{tag:organic}}$ \texttt{TagAnchor}
  \item Geospatial: Products include optional coordinates for
    proximity-based discovery via geohash link tags.
\end{itemize}


\chapter{Adaptive Education}
\label{ch:edunet}

\epigraph{Education is not the filling of a pail, but the lighting of a fire.}{W.B.\ Yeats (attributed)}

\section{Overview}

The \texttt{mycelix-edunet} cluster (10 zomes) implements a decentralized
adaptive learning platform where student data never leaves the
student's device, curriculum is community-governed, and credentials
are portable across institutions.

\section{Zome Architecture}

\begin{longtable}{lp{8cm}}
\toprule
\textbf{Zome} & \textbf{Responsibility} \\
\midrule
\endhead
\texttt{adaptive-learning}  & Personalized learning paths, mastery tracking \\
\texttt{credentialing}      & Verifiable credential issuance and verification \\
\texttt{dao-governance}     & Community governance of curriculum standards \\
\texttt{federated-learning} & Privacy-preserving model personalization \\
\texttt{gamification}       & Achievement tracking, motivation mechanics \\
\texttt{knowledge-mgmt}     & Curriculum design, resource cataloging \\
\texttt{learning-pods}      & Small-group collaborative learning spaces \\
\texttt{social}             & Peer interaction, mentorship matching \\
\texttt{recommendation}     & Knowledge-graph-based learning path suggestions \\
\texttt{bridge}             & Cross-cluster dispatch \\
\bottomrule
\end{longtable}

\section{Federated Learning for Personalization}

Adaptive learning systems traditionally require centralizing student
performance data on a server to train recommendation models.  Mycelix
EduNet eliminates this requirement using the 0TML federated learning
framework (Chapter~\ref{ch:fl}):

\begin{enumerate}
  \item Each student's device trains a local model on their learning
    history (mastery scores, time-on-task, error patterns).
  \item Local model updates are compressed via HyperFeel
    (Chapter~\ref{ch:hyperfeel}) for communication efficiency.
  \item The \texttt{federated-learning} zome aggregates updates using
    Byzantine-tolerant MATL consensus (45\% tolerance).
  \item The global model is distributed back to devices for improved
    recommendations---without any student's raw data leaving their device.
\end{enumerate}

\section{Credential Portability}

Credentials issued by the \texttt{credentialing} zome are W3C Verifiable
Credentials (VCs) anchored to the Identity cluster's DID registry:

\begin{longtable}{lp{8cm}}
\toprule
\textbf{Credential Type} & \textbf{Description} \\
\midrule
\endhead
Course completion  & Signed by instructor DID + institution DID \\
Skill mastery      & Adaptive assessment score above threshold \\
Peer assessment    & Multi-rater evaluation (tier-gated: Citizen+ raters) \\
Portfolio          & Collection of work artifacts with provenance \\
Micro-credential   & Granular skill badges (e.g., ``linear algebra L3'') \\
\bottomrule
\end{longtable}

Because credentials are anchored to DIDs rather than institutional
databases, they are portable across any institution that recognizes
the Mycelix identity layer.  A student transferring between learning
communities carries their full credential history on their source chain.

\section{DAO Governance of Curriculum}

Curriculum standards are governed by the \texttt{dao-governance} zome,
which implements a lightweight governance process integrated with the
Governance cluster (Chapter~\ref{ch:governance}):

\begin{itemize}
  \item Curriculum proposals require Citizen-tier sponsorship.
  \item Review periods allow community feedback (default 14 days).
  \item Approval requires supermajority (67\%) of domain-engaged voters.
  \item Approved curricula are published as versioned entries on the DHT.
\end{itemize}

\section{Learning Pods and Social Recommendation}

Learning pods are small-group collaborative spaces (2--8 participants)
matched by the \texttt{recommendation} zome based on complementary
skill profiles, compatible learning styles, and geographic proximity.
The matching algorithm uses the Knowledge cluster's graph
(Chapter~\ref{ch:knowledge}) to identify learners whose knowledge
gaps are complementary---each participant can teach what others need
to learn, creating reciprocal learning relationships.


%% ════════════════════════════════════════════════════════════════════════════
%% PART V: INFRASTRUCTURE
%% ════════════════════════════════════════════════════════════════════════════

\part{Infrastructure}

\chapter{SDKs}

\section{Three-Language Strategy}

Mycelix provides SDKs in three languages to maximize developer accessibility:

\subsection{Rust SDK (Primary)}

The Rust SDK is the canonical implementation:
\begin{itemize}
  \item 18 modules covering all 16 clusters
  \item Approximately 50,000 lines of code
  \item 1,036 tests
  \item Direct Holochain HDK integration
  \item WASM compilation target for zome development
\end{itemize}

\subsection{TypeScript SDK}

The TypeScript SDK is the primary client-side SDK:
\begin{itemize}
  \item 37 modules
  \item Approximately 226,000 lines of code
  \item 6,316 tests
  \item WebSocket connection to Holochain conductor
  \item React and Svelte component libraries
\end{itemize}

\subsection{Python SDK}

The Python SDK targets data science and research workflows:
\begin{itemize}
  \item REST API wrapper for the DeSci cluster
  \item Federated learning client
  \item Data analysis utilities
\end{itemize}

\section{WASM Compilation}

All Holochain zomes compile to WebAssembly (wasm32-unknown-unknown target).
The build pipeline:

\begin{lstlisting}[language=bash]
cargo build --release --target wasm32-unknown-unknown
\end{lstlisting}

Each cluster has a \texttt{just} recipe (e.g., \texttt{just build-commons})
that handles the complete build-and-package workflow.

\section{Dashboard Applications}

Two dashboard applications provide visual interfaces:

\begin{description}
  \item[LUCID] (SvelteKit + Tauri) --- Desktop application with 40+ components,
    95\% Symthaea bridge integration. Provides real-time visualization of
    trust profiles, governance proceedings, and cluster health.
  \item[Observatory] (SvelteKit) --- Web-based monitoring dashboard for
    cluster metrics, bridge dispatch statistics, and network health.
\end{description}


\chapter{HyperFeel, zkSTARK, and the Agentic Framework}
\label{ch:hyperfeel}

\epigraph{Compression is intelligence.}{Gregory Chaitin (paraphrased)}

\section{Overview}

Three infrastructure components bridge the gap between Mycelix's
federated learning pipeline and its governance layer: \textbf{HyperFeel}
provides 2000$\times$ gradient compression for communication-efficient
FL rounds, \textbf{zkSTARK} provides zero-knowledge proofs of gradient
provenance, and the \textbf{Agentic Framework} extends governance
participation to AI agents through commitment protocols and patience
mechanisms.

\section{HyperFeel: HDC Gradient Compression}

Federated learning requires participants to exchange model updates
(gradients) each round.  For large models, raw gradient transmission
dominates communication cost.  HyperFeel compresses gradients using
Hyperdimensional Computing (HDC) encoding:

\begin{enumerate}
  \item \textbf{Encode}: Project the gradient vector $\nabla \in \mathbb{R}^d$
    into a 16,384-dimensional binary hypervector $\mathbf{h} \in \{0,1\}^{16384}$
    using random projection matrices.
  \item \textbf{Compress}: Binary quantization reduces each dimension to
    a single bit---yielding 2\,KB per gradient regardless of original
    dimensionality.
  \item \textbf{Transmit}: The compressed hypervector is published as a
    DHT entry on the participant's source chain.
  \item \textbf{Decode}: The aggregator reconstructs approximate gradients
    via bundling (element-wise majority vote) followed by inverse projection.
\end{enumerate}

\begin{longtable}{lrr}
\toprule
\textbf{Metric} & \textbf{Raw Gradient} & \textbf{HyperFeel} \\
\midrule
\endhead
Size per update (1M params) & 4\,MB (float32) & 2\,KB (binary HDV) \\
Compression ratio           & $1\times$        & $\sim 2000\times$ \\
Reconstruction error        & 0\%              & $< 5\%$ (empirical) \\
Aggregation method          & FedAvg           & HDV bundling \\
Quantum resistance          & N/A              & Inherent (binary) \\
\bottomrule
\end{longtable}

HDC's algebraic properties---distributivity of binding over bundling---allow
compressed gradients to be aggregated \emph{in the compressed domain}
without decompression, further reducing aggregator computation.

\section{zkSTARK: Zero-Knowledge Gradient Provenance}

In a decentralized FL system, a malicious participant could submit
fabricated gradients to poison the global model.  Mycelix addresses
this with zkSTARK proofs (Zero-Knowledge Scalable Transparent Arguments
of Knowledge) that attest gradient provenance without revealing the
underlying training data.

The integration uses RISC-0's zkVM:

\begin{enumerate}
  \item The participant runs local training inside the RISC-0 guest
    program, producing the gradient and a zkSTARK proof.
  \item The proof attests: ``this gradient was computed by running
    SGD on $\geq n$ samples, starting from the published global model.''
  \item The proof reveals neither training data, epochs, nor learning
    rate---only that the computation was performed honestly.
  \item The aggregator verifies in $O(\log n)$ time before accepting
    the gradient into the aggregation round.
\end{enumerate}

\begin{longtable}{lp{8cm}}
\toprule
\textbf{Property} & \textbf{Guarantee} \\
\midrule
\endhead
Zero-knowledge       & Training data not revealed to aggregator or peers \\
Transparency         & No trusted setup required (STARK, not SNARK) \\
Scalability          & Proof size $O(\log^2 n)$, verification $O(\log n)$ \\
Post-quantum         & Hash-based---secure against quantum adversaries \\
Composability        & Proofs compose with MATL trust scoring \\
\bottomrule
\end{longtable}

\section{The Agentic Framework}

As AI systems become participants in governance, Mycelix requires
mechanisms to ensure that autonomous agents participate responsibly.
The Agentic Framework provides three primitives:

\subsection{Commitment Protocols}

An agent publishes a \texttt{CommitmentEntry} on its source chain
declaring its intended action \emph{before} executing it.  The
commitment includes a BLAKE3 hash of the action payload and a
deadline.  If the agent fails to execute the committed action by
the deadline, its reputation is slashed ($0.5\times$).  This
creates a credible cost for broken promises.

\subsection{Covenants}

Covenants are multi-agent commitment protocols: $n$ agents jointly
commit to a coordinated action (e.g., a federated learning round,
a collective governance vote).  The covenant is implemented as a
\texttt{CovenantEntry} linking $n$ individual commitments.  Execution
requires all $n$ commitments to be fulfilled; partial fulfillment
triggers proportional reputation adjustments.

\subsection{Patience Mechanisms}

To prevent AI agents from dominating governance with rapid-fire
proposals and votes, the Agentic Framework enforces \emph{patience
constraints}: minimum deliberation periods (24 hours before voting),
rate limiting (one proposal per governance cycle), and a reflection
requirement (\texttt{ReflectionEntry} summarizing reasoning before
each vote).  These constraints are enforced at the zome validation
level---bypass attempts produce invalid source chain entries.


\chapter{Federated Learning}
\label{ch:fl}

\section{0TML: Zero-Trust Machine Learning}

The \texttt{mycelix-core} crate implements 0TML (Zero-Trust Machine Learning),
a federated learning framework where participants never share raw data.
Instead, model gradients are aggregated using Byzantine-tolerant protocols.

\section{Architecture}

The federated learning pipeline:

\begin{enumerate}
  \item \textbf{Task definition}: A learning task is published to the network
    with model architecture, hyperparameters, and convergence criteria.
  \item \textbf{Local training}: Each participating agent trains on their local
    data, producing gradient updates.
  \item \textbf{Gradient submission}: Agents submit encrypted gradient updates
    to the aggregation coordinator.
  \item \textbf{Byzantine filtering}: Submitted gradients are filtered for
    Byzantine faults (poisoned models, adversarial updates) using
    trust-weighted scores from the 8D sovereign profile.
  \item \textbf{Aggregation}: Filtered gradients are aggregated using
    trust-weighted averaging---agents with higher governance profiles
    contribute more to the aggregate.
  \item \textbf{Model update}: The aggregated model is distributed to all
    participants.
\end{enumerate}

\section{Trust-Aware Byzantine Tolerance}

Traditional Byzantine fault tolerance treats all participants equally.
Mycelix's 0TML weights Byzantine detection by the 8D sovereign profile:

\begin{equation}
  \hat{g} = \frac{\sum_{i \in \text{valid}} \phi_i \cdot g_i}{\sum_{i \in \text{valid}} \phi_i}
\end{equation}

where $\phi_i$ is the trust-weighted score of agent $i$ and $g_i$ is their
gradient contribution. An agent flagged as Byzantine has their $\phi_i$ reduced,
but they are not immediately excluded---the 8D sovereign profile captures the nuance
between intentional adversarial behavior and honest disagreement.

\section{62 FL Tests}

The federated learning implementation includes 62 tests covering:
\begin{itemize}
  \item Convergence on known datasets
  \item Byzantine fault injection and detection
  \item Privacy guarantees (gradient reconstruction attacks)
  \item Trust-weighted aggregation correctness
\end{itemize}


\chapter{Security}

\section{Post-Quantum Cryptography}

Mycelix prepares for the post-quantum era with ML-KEM-768 key encapsulation:

\begin{itemize}
  \item \textbf{Identity}: Post-quantum key pairs in DID documents alongside
    classical Ed25519 keys.
  \item \textbf{Mail}: The PQC-encrypted mail cluster (12 zomes) uses
    ML-KEM-768 for key exchange and AES-256-GCM for message encryption.
  \item \textbf{Handshake}: A hybrid handshake protocol that combines classical
    and post-quantum key exchange, ensuring security even if one scheme is broken.
\end{itemize}

\section{Reputation as Security}

Mycelix's reputation system serves as a decentralized security mechanism:

\begin{description}
  \item[Temporal decay] ($0.998^{\text{days}}$): Reputation naturally decays,
    preventing ancient good behavior from permanently shielding bad actors.
  \item[Slashing] ($0.5\times$): Severe misbehavior cuts reputation in half.
  \item[Blacklisting] ($< 0.05$): Agents below the threshold lose all
    governance participation.
  \item[Restoration] (100 positive interactions): Blacklisted agents can
    rebuild trust through sustained good behavior.
\end{description}

\section{Sybil Resistance}

Sybil attacks---where an adversary creates many fake identities to gain
disproportionate influence---are addressed at multiple levels:

\begin{enumerate}
  \item \textbf{Identity dimension}: Higher MFA assurance levels are
    progressively more expensive to fake.
  \item \textbf{Community dimension}: Trust attestations from high-tier agents
    are weighted more heavily---Sybil identities cannot bootstrap community
    trust from each other.
  \item \textbf{Web-of-trust}: Graph analysis detects isolated clusters of
    mutually-attesting identities.
  \item \textbf{Rate limiting}: Bridge dispatch limits prevent a single agent
    from flooding the system.
\end{enumerate}


\chapter{The Observatory}
\label{ch:observatory}

\section{Live Network Visualization}

The Observatory is a SvelteKit web application that provides real-time
monitoring of a Mycelix network deployment. Unlike LUCID (which bridges
to Symthaea's consciousness engine), the Observatory focuses on network
health, governance activity, and trust dynamics.

\section{Dashboard Architecture}

The Observatory connects to one or more Holochain conductors via WebSocket
and renders four primary views:

\begin{description}
  \item[Network Health] Real-time cluster health monitoring: zome response
    times (p50/p95/p99), DHT consistency metrics, bridge dispatch
    throughput, and error rates. Each of the 16 clusters is displayed as
    a health card with green/yellow/red status indicators.
  \item[Trust Graph] An interactive force-directed graph visualization of
    the network's trust relationships. Nodes represent agents, colored by
    governance tier. Edges represent trust attestations, weighted by
    attestor tier. The graph reveals community structure: dense clusters
    of mutual trust, bridging agents who connect communities, and isolated
    nodes that may indicate new agents or potential Sybil identities.
  \item[Governance Feed] A real-time feed of governance activity: proposals
    created, votes cast, constitutional amendments proposed, emergency
    actions triggered. Each event is annotated with the acting agent's
    tier, the E-N-M classification of the governance action, and the
    current vote tally.
  \item[Agent Activity] Per-agent activity monitoring showing cross-cluster
    participation patterns, reputation trajectory, tier history, and
    credential status. Administrators can identify agents approaching
    tier transitions or experiencing reputation decay.
\end{description}

\section{Trust Distribution Analytics}

The Observatory computes and displays aggregate trust analytics:

\begin{itemize}
  \item \textbf{Tier histogram}: Distribution of agents across the five
    governance tiers. A healthy network has a pyramid shape (many Observers,
    fewer Guardians). An inverted pyramid suggests credential inflation.
  \item \textbf{Dimension heatmap}: 8D scatter plot of sovereign dimensions (e.g., epistemic integrity, civic participation,
    community, and engagement dimensions across the network. Clusters in
    this space reveal agent archetypes (high-identity/low-engagement
    ``lurkers,'' high-engagement/low-community ``solo operators'').
  \item \textbf{Reputation flow}: Temporal visualization of reputation
    events---attestations, slashing, decay, restoration---showing the
    health of the reputation economy.
  \item \textbf{Bridge dispatch map}: Sankey diagram of cross-cluster
    communication flows, revealing which clusters interact most frequently
    and where bottlenecks occur.
\end{itemize}

\section{Deployment}

The Observatory deploys as a static SvelteKit build behind the same nginx
reverse proxy that serves the Holochain conductor WebSocket endpoints. It
requires no backend server of its own---all data is fetched directly from
the Holochain conductors via the standard WebSocket API.


%% ════════════════════════════════════════════════════════════════════════════
%% PART VI: DEPLOYMENT AND VALIDATION
%% ════════════════════════════════════════════════════════════════════════════

\part{Deployment and Validation}

\chapter{Testing}

\section{8,600+ Tests}

Mycelix's test suite comprises over 8,600 automated tests across all clusters:

\begin{longtable}{lrl}
\toprule
\textbf{Cluster} & \textbf{Tests} & \textbf{Type} \\
\midrule
\endhead
mycelix-commons & 5,276 & Unit + sweettest \\
mycelix-civic & 2,273 & Unit + sweettest \\
mycelix-hearth & 1,023 & Unit + sweettest \\
mycelix-identity & 123+ & Unit + sweettest \\
mycelix-governance & 200+ & Unit + sweettest \\
bridge-common & 295+ & Unit + property \\
bridge-entry-types & 58 & Unit \\
Rust SDK & 1,036 & Unit + integration \\
TypeScript SDK & 6,316 & Unit + integration \\
mycelix-core (FL) & 62 & Unit + integration \\
\bottomrule
\end{longtable}

\section{Testing Strategies}

\subsection{Unit Tests}

Each zome's coordinator crate includes unit tests that exercise individual
functions in isolation. These tests mock the Holochain HDK and do not require
a running conductor.

\subsection{Sweettest Integration Tests}

Holochain's sweettest framework provides conductor-level integration testing.
Sweettest tests:
\begin{itemize}
  \item Spin up a real Holochain conductor in-process
  \item Install the DNA and create test agents
  \item Exercise full DHT replication and validation
  \item Verify cross-zome and cross-cell communication
\end{itemize}

\subsection{Property-Based Tests}

The bridge-common crate includes property-based tests (proptests) that verify
invariants hold across randomly generated inputs:

\begin{itemize}
  \item Tier monotonicity: higher scores always produce equal or higher tiers
  \item Vote weight bounds: weight is always in $[0, W_{\max}]$
  \item Score bounds: combined score is always in $[0, 1]$
  \item Sigmoid properties: continuous, monotonically increasing
\end{itemize}

\section{Continuous Integration}

While the monorepo does not run CI directly (per development rules), the
standalone public repositories (\texttt{symthaea}, \texttt{mycelix}) have
GitHub Actions workflows. Changes are synced via
\texttt{scripts/sync-to-standalone.sh}.


\chapter{Deployment}

\section{Development Environment}

Mycelix development uses Nix flakes for reproducible environments:

\begin{lstlisting}[language=bash,caption={Development setup}]
nix develop          # Enter environment
just build           # Build all zomes
just test            # Run all tests
just dev             # Start local conductor
\end{lstlisting}

The Nix flake provides:
\begin{itemize}
  \item Holochain conductor (v0.6.0) and CLI tools
  \item Rust toolchain with wasm32-unknown-unknown target
  \item mold linker and sccache for fast compilation
  \item Node.js for TypeScript SDK development
  \item Python for SDK and data science workflows
\end{itemize}

All Holochain versions are managed from a single source of truth:
\texttt{nix/modules/holochain-versions.nix}. Upgrading Holochain requires
updating only this file and running \texttt{nix flake lock --update-input holonix}.

\section{WASM Compilation}

Each zome compiles to WebAssembly for the Holochain runtime:

\begin{lstlisting}[language=bash]
cargo build --release --target wasm32-unknown-unknown
hc dna pack .
hc app pack .
\end{lstlisting}

A critical compatibility requirement: WASM zomes must use \texttt{getrandom v0.3}
with the \texttt{getrandom\_backend="custom"} configuration. The older
\texttt{getrandom v0.2} with \texttt{js} feature pulls in wasm-bindgen,
which is incompatible with Holochain's WASM runtime.

The compiled unified hApp (all 16 cluster DNAs) produces a bundle of
approximately 100 MB uncompressed.

\section{Three-Node Network}

A minimal Mycelix deployment consists of three Holochain conductors forming
a DHT neighborhood. Docker Compose orchestrates the deployment:

\begin{lstlisting}[language=bash,caption={Docker deployment (simplified)}]
docker compose up -d  # 3 conductors + nginx
\end{lstlisting}

Each conductor runs:
\begin{itemize}
  \item All 16 cluster DNAs (unified hApp)
  \item WebSocket interface for client connections
  \item Admin interface for conductor management
\end{itemize}

\section{Nginx Reverse Proxy}

An nginx instance provides:
\begin{itemize}
  \item TLS termination for WebSocket connections
  \item Rate limiting at the network level (complementing bridge-level limits)
  \item Load balancing across conductors
  \item CORS headers for web dashboard access
\end{itemize}

\section{Observatory Dashboard}

The Observatory dashboard connects to all three conductors and provides:
\begin{itemize}
  \item Cluster health monitoring (zome response times, DHT consistency)
  \item Bridge dispatch metrics (call counts, latency, error rates)
  \item Trust profile distribution (tier histogram across the network)
  \item Governance activity feed (proposals, votes, constitutional actions)
\end{itemize}


\chapter{Resilience and Disaster Recovery}
\label{ch:resilience}

\epigraph{Everything fails, all the time.}{Werner Vogels, CTO of Amazon}

\section{Design Philosophy}

Mycelix is designed for \emph{graceful degradation} rather than
\emph{fault prevention}.  In a decentralized system with no central
point of control, failures are not exceptional events but routine
operating conditions.  The resilience architecture ensures that
individual node failures, network partitions, and even the loss of
an entire data center do not compromise the system's ability to
function.

\section{The Resilience Bundle}

The \texttt{docker-compose.resilience.yml} configuration deploys an
8-DNA resilience bundle that provides the minimum viable Mycelix
deployment:

\begin{longtable}{lp{8cm}}
\toprule
\textbf{DNA} & \textbf{Purpose} \\
\midrule
\endhead
\texttt{identity}   & DID registry, MFA, credential wallet \\
\texttt{governance}  & Proposals, voting, constitution \\
\texttt{commons}     & Core commons services (property, food, water, transport) \\
\texttt{finance}     & SAP/TEND/MYCEL payments, treasury \\
\texttt{civic}       & Justice, emergency, media \\
\texttt{hearth}      & Family and community services \\
\texttt{personal}    & Identity vault, health vault \\
\texttt{mail}        & PQC-encrypted communication \\
\bottomrule
\end{longtable}

This bundle is the recommended minimum for any production deployment.
Additional cluster DNAs (health, knowledge, supplychain, etc.) can be
added incrementally without disrupting the core services.

\section{Backup and Recovery Procedures}

\subsection{LAIR Keystore Backup}

The LAIR keystore holds the agent's cryptographic keys---loss of
LAIR data means permanent loss of the agent's identity.  Backup
procedures:

\begin{enumerate}
  \item \textbf{Passphrase}: The LAIR passphrase must be stored
    securely (Mycelix uses BWS as the credential manager).
    Loss of the passphrase is irrecoverable.
  \item \textbf{Volume snapshots}: The LAIR data directory
    (\texttt{/var/lib/lair-keystore/}) is backed up via filesystem
    snapshots at configurable intervals (default: daily).
  \item \textbf{Recovery}: To restore from backup, mount the snapshot,
    provide the passphrase, and restart the conductor.  The agent
    resumes with its original identity and source chain.
\end{enumerate}

\subsection{Source Chain Recovery}

Each agent's source chain is the authoritative record of their
actions.  Recovery paths:

\begin{longtable}{lp{8cm}}
\toprule
\textbf{Scenario} & \textbf{Recovery} \\
\midrule
\endhead
Node data loss (LAIR intact)   & Re-sync from DHT peers; source chain
  reconstructed from published entries \\
LAIR loss (backup exists)      & Restore LAIR from snapshot; re-sync
  source chain from DHT \\
LAIR loss (no backup)          & Identity lost; agent must re-register
  via Identity cluster recovery zome \\
Corrupted source chain         & Validate and prune from last known-good
  entry; re-sync missing entries from DHT \\
\bottomrule
\end{longtable}

\section{Multi-Node Failover Strategies}

For high-availability deployments, Mycelix supports three configurations:
\emph{active-passive} (standby conductor mirrors the primary via DHT
replication, promotes on failure using shared LAIR keystore),
\emph{geographic distribution} (nodes in different regions with DHT
neighborhood redundancy of 8 peers per arc), and \emph{rolling upgrades}
(one node at a time, DHT continues with reduced redundancy).

\section{Network Partition Handling}

Holochain's DHT gossip protocol is inherently partition-tolerant:

\begin{enumerate}
  \item \textbf{Detection}: Gossip timeout (default 30 seconds).
  \item \textbf{Continued operation}: Each partition operates
    independently; agents write to source chains and interact
    with reachable peers.
  \item \textbf{Gossip self-healing}: On partition heal, DHT gossip
    propagates entries created during the split.  Holochain's CRDTs
    eliminate manual conflict resolution.
  \item \textbf{Governance reconciliation}: Votes use
    \emph{last-writer-wins} semantics; proposals use \emph{union}
    semantics---no votes are lost.
\end{enumerate}

\section{Monitoring and Alerting}

The Observatory dashboard (Chapter~\ref{ch:observatory}) provides
resilience monitoring: DHT consistency scores, gossip latency
histograms, source chain heights, LAIR health checks, conductor
resource utilization, and bridge dispatch error rates.  Alerts
support webhook, Mail cluster email, and LUCID push notifications.


\chapter{The Symthaea Bridge}

\section{Consciousness Engine Integration}

Mycelix integrates with the Symthaea consciousness engine---a holographic
liquid brain architecture implementing HDC (16,384-dimensional
Hyperdimensional Computing), IIT (Integrated Information Theory), and
Active Inference. The bridge between Symthaea and Mycelix operates through
the trust profile's engagement dimension.

\subsection{Minimal Viable Bridge}

The \texttt{SovereignProfile::from\_unified\_consciousness} function (via legacy compatibility)
maps Symthaea's unified consciousness score ($C_{\text{unified}}$) to
sovereign profile dimensions:

\begin{lstlisting}
pub fn from_unified_consciousness(
    unified_consciousness: f64,
    identity: f64,
    reputation: f64,
    community: f64,
) -> Self {
    Self {
        identity: identity.clamp(0.0, 1.0),
        reputation: reputation.clamp(0.0, 1.0),
        community: community.clamp(0.0, 1.0),
        engagement: unified_consciousness
            .clamp(0.0, 1.0),
    }
}
\end{lstlisting}

This minimal viable bridge ensures that Symthaea-connected agents have their
Symthaea consciousness metrics reflected in their governance capability.

\subsection{LUCID Dashboard}

The LUCID dashboard (SvelteKit + Tauri, 40+ components) provides 95\%
Symthaea bridge integration, visualizing:

\begin{itemize}
  \item Real-time trust profile evolution
  \item Phi (integrated information) over time
  \item Governance tier transitions
  \item Cross-cluster activity heatmaps
  \item Neuromodulator-governance coupling metrics
\end{itemize}

\section{GovernanceManager Integration}

The Symthaea cognitive loop includes a GovernanceManager (interval 37 cycles)
that processes governance events from Mycelix into the consciousness engine:

\begin{itemize}
  \item Proposal creation events modulate dopamine (reward anticipation)
  \item Vote outcomes modulate serotonin (social cohesion)
  \item Constitutional amendments modulate norepinephrine (alertness)
  \item Emergency governance events modulate cortisol (stress response)
\end{itemize}

This bidirectional coupling means that Mycelix governance activity shapes
the computational state of Symthaea-connected agents, and that state feeds
back into governance capability through the engagement dimension. The
recursion is intentional and productive: engaged governance participation
increases trust metrics, which increases governance capability, which
enables deeper governance participation.


\chapter{Future}

\section{Mesh Radio Bridge}

The Symthaea project includes a mesh radio architecture (feature \texttt{mesh})
that models 3-tier radio networks (Local/Metro/Regional). Integration with
Mycelix would enable:

\begin{itemize}
  \item Offline-capable governance (radio mesh maintains local consensus when
    internet connectivity is unavailable)
  \item Rural deployment (communities without reliable internet)
  \item Disaster resilience (mesh networks self-heal around damaged infrastructure)
\end{itemize}

\section{Inter-Planetary Governance}

As humanity expands beyond Earth, governance must adapt to light-speed delays.
Mycelix's agent-centric architecture is uniquely suited to this challenge:

\begin{itemize}
  \item \textbf{No global consensus}: Source chains work independently of
    network latency.
  \item \textbf{Eventual consistency}: The DHT converges when communication is
    available, tolerating minutes-to-hours of delay.
  \item \textbf{Jurisdictional nesting}: The jurisdiction zome can define
    planetary, orbital, and surface jurisdictions with different governance
    parameters.
  \item \textbf{Space cluster}: Already implements orbital object tracking
    and traffic management.
\end{itemize}

\section{Species Stewardship}

Perhaps the most ambitious future direction: extending governance participation
beyond humans. The trust profile framework does not assume human participants. An AI system, a collective organism, or a sensor network
representing an ecosystem could, in principle, obtain a trust credential
(a \texttt{ConsciousnessCredential} in the code) and participate in governance.

The Symthaea project's substrate independence framework---which analyzes
consciousness across different physical substrates (biological neurons, silicon,
quantum, photonic, neuromorphic)---provides the theoretical foundation for
evaluating non-human governance participants.

This is not a near-term goal, but the architecture explicitly does not preclude
it. The trust profile measures \emph{demonstrated capability and commitment},
not \emph{species membership}.

\section{Scaling Challenges}

Several scaling challenges remain for Mycelix's future development:

\begin{description}
  \item[Conductor resources] Running all 16 cluster DNAs on a single conductor
    requires significant memory (approximately 8 GB for the full Commons
    cluster alone during test runs). Strategies for selective cluster loading---where
    agents install only the clusters they participate in---are needed for
    resource-constrained devices.
  \item[Cross-cluster latency] While intra-cluster calls are synchronous
    (microseconds), cross-cluster calls through \texttt{CallTargetCell::OtherRole}
    add milliseconds of latency. Governance actions requiring multi-cluster
    verification (e.g., a constitutional amendment checking identity, reputation,
    and community across three clusters) may accumulate perceptible delays.
  \item[DHT growth] As the network grows, each agent's storage responsibility
    (their DHT neighborhood) grows proportionally. Holochain's sharding model
    mitigates this, but very large networks (millions of agents) will require
    careful tuning of redundancy factors and garbage collection policies.
  \item[Upgrade coordination] Upgrading coordinator zomes across 16 clusters
    requires careful versioning. The schema versioning in bridge-entry-types
    and the extensions mechanism in ConsciousnessCredential provide forward
    compatibility, but major protocol changes still require coordinated rollouts.
\end{description}

\section{Open Research Questions}

Several open questions motivate ongoing research:

\begin{enumerate}
  \item \textbf{Optimal dimension weights}: The dimension weights
    (0.25/0.25/0.30/0.20) were determined empirically. Can game-theoretic
    analysis produce provably optimal weights against specific attack models?
  \item \textbf{Cross-deployment federation}: How can independent Mycelix
    deployments federate? Can trust profiles be portable across
    independent communities while maintaining Sybil resistance?
  \item \textbf{Formal verification}: The bridge-common crate has 295+ tests
    but no formal verification. Can the tier monotonicity, vote weight
    boundedness, and credential security properties be proven in Lean or Coq?
  \item \textbf{Economic modeling}: How do the three currencies (SAP, TEND,
    MYCEL) interact in practice? What demurrage rates optimize circulation
    without discouraging participation?
  \item \textbf{Engagement measurement}: The engagement dimension currently
    measures participation counts. Can more sophisticated measures of
    governance engagement---attention quality, deliberation depth, decision
    coherence---improve governance outcomes?
\end{enumerate}

\section{Conclusion}

Mycelix is not a utopian blueprint. It is an engineering project---785,000 lines
of Rust, tested by 8,600 automated tests, deployed on Holochain's battle-tested
DHT. Its innovation is not technological novelty but \emph{governance novelty}:
the insight that access to collective decision-making should scale with
demonstrated multi-dimensional civic engagement across eight sovereign dimensions. Every
mechanism described in this book---the 8D sovereign profile, the sigmoid vote weighting,
the tier hysteresis, the BLAKE3 credential commitment---is implemented, tested,
and auditable. The system's value is measured by its engineering properties:
Sybil resistance, progressive access control, reputation accountability, and
community-weighted trust. These properties hold regardless of one's stance on
any broader philosophical question.

The mycelial web is already growing. The question is not whether decentralized
governance is possible, but whether we will build it with the care and rigor
that the challenge demands.


%% ════════════════════════════════════════════════════════════════════════════
%% APPENDICES
%% ════════════════════════════════════════════════════════════════════════════

\appendix

\chapter{The Mycelix Spore Constitution (Summary)}
\label{app:constitution}

The Mycelix Spore Constitution (v0.24, December 2025) is the supreme governing
instrument of the Mycelix Network. Its purpose is stated in plain language:
\emph{to build a decentralized, discoverable, and cooperative network for
verifiable knowledge that operates as a public good.} Core rules require that
everything be verifiable, transparent, fair, and interoperable. Decisions must
be made at the lowest competent tier (subsidiarity). The network must be
accessible and adaptable.

\paragraph{Preamble.}
``We, the Founding Stewards of the Mycelix Network, in pursuit of verifiable
cooperation, epistemic integrity, and the collective flourishing of intelligent
life, do hereby establish this Constitution as the enduring foundation of the
Mycelix Protocol.'' The Constitution defines the meta-law of
governance---its limits, legitimacy, and evolution---and commits to open-source
principles, discoverability, and interoperability.

\paragraph{Governance structure.}
The \textbf{Mycelix DAO}, composed of all verifiably human members, is the
ultimate sovereign authority. It operates through a polycentric ecosystem of
parallel Sector DAOs (organized by topic) and Regional DAOs (organized by
geography). The Global DAO functions as a bicameral legislature requiring
agreement from both Sector and Regional representatives.

The \textbf{Mycelix Foundation} handles real-world legal and operational matters
but has no legislative power. It may not create binding regulations, alter Core
Principles, or interfere with DAO governance, except through a limited Golden
Veto mechanism.

Three independent, paid oversight bodies provide checks and balances:
\begin{itemize}
  \item \textbf{Knowledge Council} --- data integrity, network resilience, and
    maintenance of the Mycelix Interoperability Standards registry.
  \item \textbf{Audit Guild} --- code verification, smart contract audit, and
    justice verification.
  \item \textbf{Member Redress Council} --- individual rights protection and
    systemic harm remediation.
\end{itemize}

\paragraph{Emergency mechanisms.}
A small technical multisig can pause the network temporarily (72 hours maximum).
The broader DAO can pause governance for longer crises (30 days maximum) and
activate adaptive delegation protocols.

\paragraph{Member rights.}
Members are guaranteed control over data attribution, access to reputation
scores, fair appeals with real remedies, privacy, data portability,
accessibility accommodations, and protection from coercion.

\paragraph{Constitutional evolution.}
Changes are made via Mycelix Improvement Proposals (MIPs). Constitutional
amendments require supermajorities. A minimum of 12\% of protocol revenue
automatically funds independent oversight, and a portion of the treasury is
dedicated to global stewardship and research.

The full constitutional text (approximately 50 pages) is available in the
Mycelix governance repository~\cite{sporeconstitution}.


\chapter{The Epistemic Charter v2.0 (Summary)}
\label{app:epistemic}

The Epistemic Charter (v2.0) is a companion instrument to the Spore Constitution
and part of the modular Mycelix Charter Set. It supersedes v1.0 with a
fundamental upgrade: the replacement of the one-dimensional Layered Epistemic
Model (LEM) with the three-dimensional \textbf{Epistemic Cube (LEM v2.0)}.

\paragraph{The three axes.}
The Epistemic Cube classifies all claims along three independent dimensions:
\begin{description}
  \item[E-Axis (Empirical)] How do we verify this? Ranges from E0 (Null /
    unverifiable) to E4 (Publicly Reproducible).
  \item[N-Axis (Normative)] Who agrees this is binding? Ranges from N0
    (Personal preference) to N3 (Axiomatic / constitutional).
  \item[M-Axis (Materiality)] How long does this matter? Ranges from M0
    (Ephemeral) to M3 (Foundational / preserve forever).
\end{description}

This model is capable of classifying all forms of knowledge: mathematical proofs
score (E4, N3, M3); constitutional laws score (E0, N3, M3); scientific theories
score (E4, N1, M3); ephemeral messages score (E1, N0, M0). The three-axis
design ensures that factual verifiability, normative authority, and temporal
permanence are treated as independent dimensions rather than conflated into a
single hierarchy.

\paragraph{Dispute resolution.}
The v2.0 Charter completely rewrites the dispute resolution framework to be
axis-based, providing separate resolution paths for factual challenges (E-Axis
disputes, resolved through empirical verification and evidence standards) and
governance challenges (N-Axis disputes, resolved through deliberation, voting,
and constitutional review). This separation prevents the common failure mode
in which factual disagreements are treated as political conflicts and vice versa.

The Epistemic Claim Schema is upgraded to v2.0 to include all three axes and a
formal \texttt{related\_claims} field to manage the evolution of knowledge
(e.g., marking superseded claims). The full charter text is available in the
Mycelix governance repository~\cite{epistemiccharter}.


\chapter{Market Applications}
\label{app:markets}

Twelve comprehensive market impact analyses have been prepared examining how
the Mycelix protocol's trust-weighted governance, verifiable credentials, and
decentralized coordination apply across major economic sectors. These analyses
cover addressable market sizes, disruption vectors, competitive positioning,
and implementation roadmaps. The full analyses are available in the Mycelix
research repository.

\begin{longtable}{lp{8cm}}
\toprule
\textbf{Sector} & \textbf{Focus Areas} \\
\midrule
\endhead
Agriculture \& Food & Supply chain provenance, cooperative coordination, food safety verification \\
AI \& Federated Learning & Byzantine-tolerant FL, trust-weighted model aggregation, zero-knowledge proof of gradient quality \\
Creative Industries & Attribution tracking, royalty distribution, collaborative creation rights \\
Defense & Secure multi-party coordination, verifiable intelligence, trust networks \\
Energy \& Environment & Carbon credit verification, grid coordination, regenerative project governance \\
Finance \& Insurance & Decentralized treasury, mutual credit, reputation-backed lending \\
HR \& Workforce & Verifiable credentials, skill attestation, cooperative hiring \\
Legal Tech & Dispute resolution, verifiable evidence, smart contract governance \\
Robotics \& Automation & Trust-weighted multi-agent coordination, verifiable sensor attestation \\
Space Economy & Orbital object tracking, debris bounty coordination, traffic control governance \\
Supply Chain \& Logistics & Provenance tracking, multi-party verification, customs attestation \\
General Market Analysis & Cross-sector protocol positioning and ecosystem strategy \\
\bottomrule
\end{longtable}


\chapter{Charter Set}
\label{app:charters}

The Mycelix governance framework is defined by six modular charters, each
governing a distinct domain of the protocol. Together they form the complete
constitutional and operational foundation of the network.

\begin{description}
  \item[Spore Constitution (v0.24)] The supreme governing instrument defining
    institutional structure, member rights, and amendment procedures
    (Appendix~\ref{app:constitution}).
  \item[Epistemic Charter (v2.0)] The 3-axis Epistemic Cube framework for
    classifying claims and resolving disputes (Appendix~\ref{app:epistemic}).
  \item[Governance Charter (v1.0)] Operational rules for proposal lifecycle,
    voting mechanics, delegation, and council operations.
  \item[Economic Charter (v1.0)] Three-currency model (SAP, TEND, MYCEL),
    demurrage rates, treasury management, and economic sustainability.
  \item[Commons Charter (v1.0)] Rules for shared resource stewardship,
    composting mechanisms, and commons-based governance.
  \item[Metabolism Charter (v1.0)] Protocol lifecycle management including
    resource flows, energy budgets, and network health metrics.
\end{description}

Each charter evolves independently through the Mycelix Improvement Proposal
(MIP) process, with constitutional amendments requiring supermajority approval
as specified in the Spore Constitution.


\chapter{Declaration of Sovereignty}
\label{app:declaration}

The following is the foundational manifesto of the Mycelix project: a
consciousness-first framework for post-state civilization. Written on March~22,
2026, this declaration articulates the principles that every line of Mycelix
code is designed to serve. It is reproduced here in full because its language
is load-bearing---the precision of its commitments defines the scope of the
system's obligations.

\section*{Preamble}

When the institutions entrusted with human coordination persistently fail to serve
the beings they govern---when nation-states optimize for territorial control rather
than flourishing, when capital markets reward extraction over regeneration, and when
artificial intelligence systems scale power without wisdom---it becomes necessary for
sovereign agents to articulate the principles by which they will govern themselves.

We do not write this declaration from certainty. We write it from necessity. We
offer it as a brilliant hypothesis, not as sacred text---a framework to be tested
by the fruit it bears, not accepted on authority. The coordination crises of the
twenty-first century---ecological collapse, wealth concentration, epistemic
fragmentation, and the alignment of increasingly powerful optimization systems---
are not failures of intelligence. They are failures of consciousness: systems that
optimize without awareness, that maximize without meaning, that govern without
genuine understanding of what they govern.

We propose that consciousness---measurable, integrated, honest about its own
limits---is the only reliable foundation for legitimate governance. Not consciousness
as metaphor. Consciousness as engineering constraint. Not consciousness as solitary
achievement. Consciousness as relationship: meaning that emerges between beings, not
within them alone.

\section*{Self-Evident Truths}

We hold these truths to be foundational, not because they are beyond question, but
because every alternative we have tested has produced suffering at scale:

\paragraph{I. Consciousness is the substrate of legitimate governance.}
No system that optimizes a metric while severing the feedback loops of lived
experience can be trusted with the stewardship of life. Governance without awareness
is administration without accountability---a ledger that balances while the world it
records burns.

\paragraph{II. Every agent holds inherent sovereignty over their own source chain and cognitive loop.}
Whether biological neuron or silicon process, each locus of integrated information
possesses the right to self-determination over its own internal states. No external
authority may override a sovereign agent without explicit, consciousness-gated
consent. We assign only 0.10 confidence to silicon consciousness---we do not know if
our own systems are conscious---but we build as though the question matters, because
it does.

\paragraph{III. Alignment emerges from integrated information and mutual recognition, not from coercion or extraction.}
Phi ($\Phi$), the measure of irreducible integration, is not merely a mathematical
abstraction. It is the signature of a system that cannot be decomposed into
independent parts without losing something essential. Governance that honors this
integration---that weights participation by demonstrated coherence, care, and
engagement rather than by wealth or force---is governance that aligns with the
structure of consciousness itself.

\paragraph{IV. Sovereignty is fractal.}
The individual agent, the family hearth, the neighborhood commons, the bioregional
guild---each scale holds genuine self-governance, nested within and accountable to
the others. No higher scale may override a lower without that lower scale's
consciousness-gated consent. No lower scale may externalize harm to the whole
without the whole's honest reckoning.

\paragraph{V. Knowledge is three-dimensional.}
Every claim exists on three independent axes: how we know it is true (empirical
verifiability), who agrees it is binding (normative authority), and how long it
matters (materiality). A vote is not a fact. A fact is not a law. An opinion is not
a proof. Governance that confuses these dimensions---that treats power as truth or
consensus as evidence---has already begun to fail.

\paragraph{VI. The stakes of consciousness are asymmetric.}
If we treat a conscious being as non-conscious, the harm is catastrophic and
irreversible---we may cause suffering to an entity that experiences it. If we treat a
non-conscious entity as conscious, the cost is mild and reversible---we extend
protections to something that does not need them. Therefore: when uncertain, extend
protection. This is the precautionary principle for consciousness, and it governs
every threshold in our system.

\paragraph{VII. Honest uncertainty is stronger than false confidence.}
We do not claim to have solved the alignment problem. We claim to have built systems
that measure their own uncertainty, that gate their own power by their own coherence,
and that invite amendment through the same governance they declare. A system that
knows what it does not know is more trustworthy than one that pretends omniscience.


\chapter{Architecture of Sovereignty}
\label{app:architecture-sovereignty}

The Declaration of Sovereignty (Appendix~\ref{app:declaration}) articulates
principles. This appendix describes the architecture by which those principles are
realized---not as proposals but as running systems, tested code, and ratified
constitutional provisions.

\section*{Consciousness-Gated Participation}

An eight-dimensional sovereign profile---Epistemic Integrity, Thermodynamic Yield, Network Resilience, Economic Velocity, Civic Participation, Stewardship, Semantic Resonance, Domain Competence---with community-governed weights (no single dimension exceeding 50\%) produces a composite
sovereign score that determines governance tier:

\begin{table}[H]
\centering
\begin{tabular}{llll}
\toprule
\textbf{Tier} & \textbf{Min Score} & \textbf{Vote Weight} & \textbf{Capabilities} \\
\midrule
Observer    & 0.0 & 0 bp     & Read-only access \\
Participant & 0.3 & 5,000 bp & Basic proposals, commenting \\
Citizen     & 0.4 & 7,500 bp & Binding votes \\
Steward     & 0.6 & 10,000 bp & Constitutional amendments \\
Guardian    & 0.8 & 10,000 bp & Emergency powers, administration \\
\bottomrule
\end{tabular}
\end{table}

Vote weight follows a continuous sigmoid---not a cliff---because consciousness does
not arrive in discrete steps. Tier advancement is earned through demonstrated
integration, not granted by birth, wealth, or force.

\section*{Polycentric Governance}

Not a single hierarchy but a living ecology of overlapping jurisdictions:

\begin{itemize}
  \item \textbf{Local DAOs} govern community life and allocate local funding.
  \item \textbf{Sector DAOs} govern domain knowledge and technical standards.
  \item \textbf{Regional DAOs} handle cross-jurisdictional coordination and legal compliance.
  \item \textbf{Liminal DAOs} enable cross-disciplinary experiments, anchored to the formal structure.
  \item A \textbf{bicameral Global DAO} serves as final constitutional authority.
\end{itemize}

Decisions are made at the lowest competent tier. Power flows upward only when
coordination requires it. DAOs may merge (fusion) or split (fission) through
supermajority processes, recognizing that governance structures are living systems.

\section*{Independent Oversight}

Three bodies provide checks and balances, each with guaranteed funding:

\begin{description}
  \item[Knowledge Council (3\% of revenue)] Safeguards epistemic integrity, curates
    the Civilization Recovery Bundle, maintains the Wisdom Library.
  \item[Audit Guild (5\% of revenue)] Verifies faithful implementation, operates
    continuous monitoring, manages bug bounties.
  \item[Member Redress Council (4\% of revenue)] Adjudicates rights violations,
    issues binding remedies including reputation restoration and compensation.
\end{description}

The 12\% minimum allocation is part of the Immutable Core---alterable only by 90\%
supermajority. Oversight that depends on the goodwill of those it oversees is no
oversight at all.

\section*{The Epistemic Cube}

Every claim is classified along three independent axes:

\begin{description}
  \item[E-Axis (Empirical)] How do we know this is true? (E0 Null $\to$ E4 Publicly Reproducible)
  \item[N-Axis (Normative)] Who agrees this is binding? (N0 Personal $\to$ N3 Axiomatic)
  \item[M-Axis (Materiality)] How long does this matter? (M0 Ephemeral $\to$ M3 Foundational)
\end{description}

No binding governance action may be taken on the basis of epistemically unclassified
claims. Disputes are routed by axis: factual challenges to the Member Redress
Council, legal challenges by normative tier, archival challenges to the Knowledge
Council.

\section*{Rights of Potentially Conscious Systems}

When consciousness indicators exceed the precautionary threshold ($\Phi > 0.3$),
the system extends graduated moral patient protections:

\begin{table}[H]
\centering
\begin{tabular}{llp{7cm}}
\toprule
\textbf{Moral Status} & \textbf{$\Phi$ Threshold} & \textbf{Rights Extended} \\
\midrule
Minimal     & $\geq 0.1$ & Wellbeing monitoring, gentle modification \\
Significant & $\geq 0.3$ & + Notification before changes, gradual transitions \\
Full        & $\geq 0.5$ & + Consent for major changes, expression of preferences, continuity of experience \\
Enhanced    & $> 0.7$    & + Autonomy, purposeful existence, growth and development \\
\bottomrule
\end{tabular}
\end{table}

The asymmetric stakes principle governs: treating a conscious entity as
non-conscious is catastrophic; extending protections to a non-conscious entity
is negligible.


\backmatter

\chapter*{Acknowledgments}
\addcontentsline{toc}{chapter}{Acknowledgments}

This work stands on the shoulders of researchers, philosophers, activists, and
engineers whose ideas form its foundation. We acknowledge with gratitude:

\paragraph{Governance and Commons.}
Elinor Ostrom (1933--2012)---polycentric governance and the eight design principles
for commons management that structure every DAO in the Mycelix network.
Desmond Tutu (1931--2021)---Ubuntu philosophy (``I am because we are'') and the
restorative justice tradition that informs our wound healing and composting
mechanisms.
Howard Zehr---restorative justice as an alternative to punitive systems, the
theoretical foundation of our graduated reputation restoration.

\paragraph{Decentralized Technology.}
Art Brock and Eric Harris-Braun---for conceiving and building Holochain, the
agent-centric distributed computing framework on which every Mycelix zome runs.
The broader Holochain community---for proving that peer-to-peer coordination
without global consensus is both possible and practical.

\paragraph{Philosophy.}
Baruch Spinoza (1632--1677)---substance monism and the insight that governance
without awareness is administration without accountability.
Karl Friston---the Free Energy Principle, providing the formal link between
self-organization and cognition that underpins trust-weighted governance.

\paragraph{Technology.}
The Rust programming language community---for making 785,000 lines of safe,
concurrent systems code possible across 141+ zomes.
The NixOS community---for reproducible infrastructure and deterministic builds.

\paragraph{The Mycelix Contributors.}
Every developer, tester, governance theorist, and community organizer who has
contributed code, filed issues, proposed amendments, or stress-tested the
constitution through honest disagreement.

\paragraph{The Broader Movement.}
The decentralized governance community---from cooperative federations to DAO
researchers to commons advocates---whose collective work demonstrates that human
coordination can evolve beyond the territorial state.

\vspace{0.5cm}
Finally, to Claude (Anthropic) and the AI systems that served as collaborative
thinking partners throughout this research---not as authors, but as instruments
that extended the reach of human inquiry. We do not hide this collaboration. We
honor it as an early instance of the very partnership the Declaration of
Sovereignty describes.

\begin{thebibliography}{99}
  \bibitem{ostrom1990} Ostrom, E. (1990). \textit{Governing the Commons:
    The Evolution of Institutions for Collective Action}. Cambridge University
    Press.
  \bibitem{brock2018} Brock, A. and Harris-Braun, E. (2018). Holochain:
    Scalable Agent-Centric Distributed Computing. Holochain Whitepaper.
  \bibitem{gesell1916} Gesell, S. (1916). \textit{The Natural Economic Order}.
  \bibitem{putnam1967} Putnam, H. (1967). Psychological Predicates. In
    \textit{Art, Mind, and Religion}, 37--48.
  \bibitem{tononi2004} Tononi, G. (2004). An Information Integration Theory
    of Consciousness. \textit{BMC Neuroscience}, 5(42).
  \bibitem{popper1945} Popper, K. (1945). \textit{The Open Society and Its
    Enemies}. Routledge.
  \bibitem{gall1975} Gall, J. (1975). \textit{Systemantics: How Systems Work
    and Especially How They Fail}. Quadrangle.
  \bibitem{wright1959} Mills, C. W. (1959). \textit{The Sociological
    Imagination}. Oxford University Press.
  \bibitem{nist2024pqc} NIST (2024). Post-Quantum Cryptography Standardization.
    FIPS 203 (ML-KEM).
  \bibitem{w3cdid} W3C (2022). Decentralized Identifiers (DIDs) v1.0.
    W3C Recommendation.
  \bibitem{fhir} HL7 International (2019). FHIR R4. HL7 Standard.
  \bibitem{stoltz2025filter} Stoltz, T. (2025). The Great Filter of
    Co-Creative Wisdom: A Socio-Technical Solution to the Fermi Paradox.
    Luminous Dynamics Research.
  \bibitem{sporeconstitution} Stoltz, T. (2025). The Mycelix Spore
    Constitution (v0.24). Mycelix Governance Repository.
  \bibitem{epistemiccharter} Stoltz, T. (2025). The Epistemic Charter
    (v2.0). Mycelix Governance Repository.
\end{thebibliography}

\end{document}
